\section{Quantum preliminaries}  \label{sec:prelims}
\subsection{Quantum measurements}
The most general notion of a quantum measurement is a POVM
measurement, which consists of a set of Hermitian operators $\{M_a\}_{a \in S}$
indexed by outcomes $a$ from a set $S$. These satisfy the
conditions
\[ \forall a,\; M_a \succeq 0, \qquad \sum_a M_a = I. \]
To refer to the measurement as a whole
we will use the letter $M$, without the subscript indicating the
outcome. For a state $\ket{\psi}$, the
probability that the measurement $M$ returns outcome $a$ is
\[ \Pr[\ba] = \bra{\psi} M_{\ba} \ket{\psi}. \]
A POVM is said to be \emph{projective} if each element $M_a$ is an
orthogonal projector, i.e. $M_a^2 = M_a$. Note that this implies that
$M_a M_b = 0$ for any $a \neq b$, i.e. that the projectors are
pairwise orthogonal. Naimark's theorem says that any POVM measurement
can be simulated by a projective measurement on an enlarged space.
\begin{theorem}[Naimark]\label{thm:naimark}
  Suppose $\{M_a\}$ is a POVM acting on a Hilbert space $\calH$. Then
  there exists a projective measurement $\{M'_a\}$ acting on the space
  $\calH \ot \calH_{\rmaux}$ together with a state $\ket{\rmaux}$ such
  that for all states $\ket{\psi} \in \calH$ and all outcomes $a$, the
  post-measurement state after applying $M$ and $M'$ is the same:
  \begin{equation} \sqrt{M_a} \ket{\psi}\bra{\psi} \sqrt{M_a} = \tr_{\reg{aux}}(M'_a (\ket{\psi}\bra{\psi}
    \ot \ket{\rmaux}\bra{\rmaux})M'_a). \label{eq:naimark}\end{equation}
  As a consequence, $M$ and $N$ induce the same distribution over
  outcome probabilities:
  \[ \bra{\psi} M_a \ket{\psi} = (\bra{\psi} \ot \bra{\rmaux}) N_a
    (\ket{\psi} \ot \ket{\rmaux}). \]
  Moreover, given any upper-bound $n$ on the number of outcomes of
  $M_a$, there is a universal choice of the state $\ket{\rmaux}$ that
  works for all POVMs $M_a$ with at most $n$ outcomes. The projective measurement
  $M'_a$ and state $\ket{\rmaux}$ together constitute a \emph{Naimark
    dilation} of the POVM $M_a$.
\end{theorem}

  \begin{theorem}[Partial Naimark]\label{thm:partial-naimark}
    Suppose $\{M_{a_1, a_2}\}$ is a POVM acting on a tensor product
  Hilbert space $\calH = \calH_1 \ot \calH_2$ of the form $M_{a_1,
    a_2} = \Pi_{a_1} \ot A^{a_1}_{a_2}$, where the operators $\{\Pi_{a_1}\}_{a_1}$ is a
  projective measurement. Then there is a Naimark dilation $M'_{a_1,
    a_2}$ and a state $\ket{\rmaux}$ as above, with the property that
  $M'_{a_1, a_2} = \Pi_{a_1} \ot A'^{a_1}_{a_2}$ for projectors
  $A'^{a_1}_{a_2}$ acting on $\calH_2 \ot \calH_{\rmaux}$.
\end{theorem}
\begin{proof}
  The condition that $\{\Pi_{a_1}\}$ forms a projective measurement
  implies that for each $a_1$, $\{A^{a_1}_{a_2}\}$ is a
  POVM. By~\Cref{thm:naimark}, for each $a_1$ there exists a POVM
  $A'^{a_1}_{a_2}$ dilating $A^{a_1}_{a_2}$, and all of these POVMs
  act on the same universal auxiliary state $\ket{\rmaux}$. Now,
  define $M'_{a_1, a_1} = \Pi_{a_1} \ot A'^{a_1}_{a_2}$. For every
  state $\ket{\psi}$, we have that
  \begin{align*}
    \tr_{\reg{aux}}(M'_{a_1, a_2} (\ket{\psi}\bra{\psi}) M'_{a_1,
    a_2}) &= \tr_{\reg{aux}}( A'^{a_1}_{a_2} (\Pi_{a_1} \ket{\psi}
            \bra{\psi} \Pi_{a_1} \ot \ket{\rmaux} \bra{\rmaux}) A'^{a_1}_{a_2}) \\
          &= \sqrt{A^{a_1}_{a_2}} (\Pi_{a_1} \ket{\psi} \bra{\psi}
            \Pi_{a_1}) \sqrt{A^{a_1}_{a_2}} \\
          &= \sqrt{M_{a_1, a_2}} \ket{\psi} \bra{\psi} \sqrt{M_{a_1,
            a_2}},
  \end{align*}
  where in going from the first to the second line we used~\Cref{thm:naimark}.
\end{proof}


For the purposes of this paper, we will need to specialize the POVM
notation introduced above in several ways. First, we will often work
with families of POVM measurements indexed by questions in an
interactive proof protocol. These will be denoted $\{M^q_a\}$, where
$q$ indexes the question and $a$ the outcome.
(We note that the reverse convention ``$\{M_q^a\}$", which we will \emph{not} use, is also common in the literature.)
In many cases, the outcomes will consist of
tuples of elements, some of which we may wish to discard. We use the
convention that if an outcome element is not written, it
is understood to be \emph{summed} over. Thus, if $\{M^{x}_{a,b}\}$
is a family of POVMs, we would have
\begin{equation}\label{eq:spider-man}
M^{x}_{a} := \sum_b M^{x}_{a,b}.
\end{equation}

\begin{notation}\label{not:marginalize}
We will also often consider situations where some of the information
in a measurement outcome is discarded. In particular, given a
POVM $\{M_f\}$ whose outcomes are functions $f:U \to V$ over some
domain $U$, and given a point $x \in U$, we will denote by $\{M_{f(x)=
  y}\}_{y \in V}$ the measurement corresponding to applying $M$ to
obtain a function $f$, and returning the value of $f$ at
$x$. Formally, the POVM elements of this measurement are given by
\[ M_{[f(x) = a]} = \sum_{f: f(x) = a} M_f. \]
(We note that \Cref{eq:spider-man} can be viewed as a special case
in which the ``discarding function"~$f$ simply removes the second coordinate.
For this case, it is simpler to use the convenient notation in \Cref{eq:spider-man}
than the more cumbersome bracket notation given here.)
\end{notation}

The following lemma contains a useful fact about marginalized
projective measurements.
\begin{lemma}\label{lem:marginal-tensor-product}
  Let $M_{a,b}$ be a projective measurement on a tensor product
  Hilbert space $\calH_1 \ot \calH_2$, and suppose that for all $a$, $M_{a} = A_a
  \ot B_a$ where $A_a$ is a rank-one matrix on $\calH_1$. Then for all
  $a, b$, $M_{a,b} = A_a \ot C_{a,b}$ with $C_{a,b}$ projectors.
\end{lemma}
\begin{proof}
  By the Schmidt decomposition, we can write $M_{a,b}$ as
  \begin{align*}
    M_{a,b} &= \sum_j \ket{\psi_{a,b,j}}\bra{\psi_{a,b,j}} \\
    &= \sum_{j,k}  \sigma_{jk}              \ket{u_{a,b,j,k}}\bra{u_{a,b,j,k}}
      \ot\ket{v_{a,b,j,k}}\bra{v_{a,b,jk}}.
  \end{align*}
  Write the rank-one matrix $A_a$ as an outer product
  $\ket{\psi_a}\bra{\psi_a}$. Then we have
  \[ \sum_b M_{a,b} = \sum_{j,k, b} \sigma_{jk}              \ket{u_{a,b,j,k}}\bra{u_{a,b,j,k}}
      \ot\ket{v_{a,b,j,k}}\bra{v_{a,b,jk}} = 
      \ket{\psi_a}\bra{\psi_a} \ot B_a. \]
    Taking the partial trace on the $B$ system, we have
    \[  \tr_B(\sum_b M_{a,b}) = \sum_{j,k,b} \sigma_{jk}
      \ket{u_{a,b,j,k}}\bra{u_{a,b,j,k}} =
      \ket{\psi_a}\bra{\psi_a}. \]
    Suppose we multiply on the left by $\bra{v}$ and on the right by
    $\ket{v}$, for $\ket{v}$ orthogonal to $\ket{\psi_a}$. Then the
    RHS is 0 while the LHS is a sum $\sum_{j,k,b} \sigma_{jk}
    |\braket{v |u_{a,b,j,k}}|^2$ of nonnegative terms. Hence, each of
    these terms must be zero. Thus, the equation can only hold if all
    the vectors $\ket{u_{a,b,j,k}}$ are multiples of
    $\ket{\psi_a}$. This implies that
    \[ M_{a,b} = \ket{\psi_a} \bra{\psi_a} \ot C_{a,b}\]
    for some $C_{a,b}$ as desired.
\end{proof}

\subsection{Nonlocal games and $\MIP^*$}

Now, we augment \Cref{def:two-player-one-round,def:mip-protocols}
to allow for provers to share quantum resources.

\begin{definition}
Given a game~$\game$, a \emph{quantum strategy}
is one in which Alice and Bob are allowed to share entanglement but not to communicate.
We can model their behavior with the \emph{strategy} $\calS= (\rho, A, B)$.
Here,
\begin{itemize}
\itemsep -.5pt
\item[$\circ$] Write $\mathcal{H}_A$ for Alice's local Hilbert space and $\mathcal{H}_B$ for Bob's.
		Then $\rho$ is a (possibly entangled) state in $\mathcal{L}(\mathcal{H}_A \otimes \mathcal{H}_B)$.
\item[$\circ$] The set $A$ contains a matrix $A_a^x$ for each question~$x$ and answer~$a$,
		with the guarantee that for each question~$x$, $A^x := \{A_a^x\}_a$ is a POVM. (Likewise for~$B$.)
\end{itemize} 
Alice and Bob perform their strategy as follows:
given question~$x$, Alice performs the POVM $\{A_a^x\}_a$ and returns her measurement outcome to the verifier.
Bob plays similarly.
The \emph{value} of their strategy, denoted $\valstrat{\game}{\calS}$, is the probability that they pass the test,
over the randomness in~$\game$ and in their measurement outcomes.
\begin{align*}
\valstrat{\game}{\calS}
&= \E_{(\bx_0, \bx_1) \sim \mathrm{Alg}_{\mathrm{Q}}} \Pr_{\ba_0, \ba_1}[\mathrm{Alg}_{\mathrm{A}}(\bx_0, \bx_1, \ba_0, \ba_1) = 1]\\
&= \E_{(\bx_0, \bx_1) \sim \mathrm{Alg}_{\mathrm{Q}}} \sum_{\substack{a_0, a_1,\\\mathrm{Alg}_{\mathrm{A}}(\bx_0, \bx_1, a_0, a_1) = 1}}
	\tr(A_{a_0}^{\bx_0} \otimes A_{a_1}^{\bx_1} \cdot \rho),
\end{align*}
where in the first line, $(\ba_0, \ba_1)$ is the distribution on answers given questions~$\bx_0, \bx_1$.
We write $\val{\game}$ for the infimum of $\valstrat{\game}{\calS}$ over all strategies~$\calS$.
We define value analogously for interactive proofs.

We say that $L \in \MIP_{c, s}^*$ if there is an quantum interactive proof~$\game$ that decides it.
This means that the following three conditions are true.
\begin{itemize}
\item[$\circ$] (Completeness) Suppose $\mathsf{input} \in L$.  Then there is a quantum strategy for~$\game$ with value at least~$c$.
\item[$\circ$] (Soundness) Suppose $\mathsf{input} \notin L$.  Then every quantum strategy for~$\game$ has value at most~$s$.
\item[$\circ$] All of $\qlength{\game}$, $\alength{\game}$, $\qtime{\game}$, and $\atime{\game}$ are $\poly(n)$.
\end{itemize}
If $c - s$ is a constant, then we will suppress the dependence on them and just say that $L \in \MIP^*$.
\end{definition}

\begin{remark}
A game~$\game$ is \emph{symmetric} if its distribution on questions treats Alice and Bob symmetrically.
In this case, we may assume without loss of generality that Alice and Bob's strategies are also symmetric,
i.e.\ that $A_a^x = B_a^x$ for all questions~$x$ and answers~$a$.
This allows us to represent their measurements by a single set of matrices $M$ (for which $M_a^x = A_a^x = B_a^x$).
As a further simplification, by applying Naimark's dilation theorem to Alice and Bob's strategy 
we can assume that their shared state~$\psi$ is pure and their measurements are projectors.
\end{remark}

Occasionally, it will be useful to speak of the distribution over
measurement outcomes induced by a strategy independently of any
particular game. For this, we introduce the notion of a bipartite correlation
\begin{definition}
  Given a strategy $\calS = (\rho, A, B)$, the bipartite correlation produced by it is
  the function $P(a, b | x, y) = \tr(A^x_a \ot B^x_y \cdot \rho)$.
\end{definition}
If two strategies produce the same bipartite correlation, they have
the same value for any game they are used for. Naimark's theorem
(\Cref{thm:naimark}) implies for any strategy $\calS$, there exists a
strategy $\calS'$ using only projective measurements that produces the
same correlation:

\begin{corollary}[Naimark's theorem for strategies]\label{cor:bipartite-naimark}
  Suppose $\{M_a\}$ and $\{N_b\}$ are two POVMs acting on the $A$ and
  $B$ factors of a tensor Hilbert space $\calH_{A} \ot \calH_{B}$,
  respectively. Then for any Naimark dilation of $\{M_a\}$ given by
  projectors $M'_a$ and an auxiliary state $\ket{\rmaux_A} \in
  \calH_{\rmaux_A}$, and any Naimark dilation of $\{N_b\}$ given by
  projectors $\N'_b$ and an auxiliary state
  $\ket{\rmaux_B} \in \calH_{\rmaux_B}$, it holds that for any bipartite state $\ket{\psi} \in
  \calH_{A} \ot \calH_{b}$, the post-measurement state after applying
  $M \ot N$ to $\ket{\psi}$ and $M' \ot N'$ to $\ket{\psi} \ot
  \ket{\rmaux_A} \ot \ket{\rmaux_B}$ is the same:
  \[ \sqrt{M_a} \ot \sqrt{N_b}  \ket{\psi} \bra{\psi}  \sqrt{M_a} \ot
    \sqrt{N_b} = \tr_{\rmaux}[ (M'_a \ot N'_b) (\ket{\psi}\bra{\psi}
    \ot \ket{\rmaux_{A}}\bra{\rmaux_{A}} \ot
    \ket{\rmaux_B}\bra{\rmaux_B} )(M'_a \ot N'_b)]. \]
  Moreover, such dilations exist by \Cref{thm:naimark}.
  As a consequence, $M, N$ and $M', N'$ induce the same joint
  distribution over outcome probabilities:
  \[ \bra{\psi} M_a \ot N_b \ket{\psi} = (\bra{\psi} \ot
    \ket{\rmaux_A} \ot \ket{\rmaux_B}) M'_a \ot N'_b (\ket{\psi} \ot
    \ket{\rmaux_A} \ot \ket{\rmaux_B}). \]
\end{corollary}
\begin{proof}
  The existence of such dilations follows immediately
  from~\Cref{thm:naimark}. To deduce the equality of post-measurement
  states, we apply~\Cref{eq:naimark} twice, and use the fact that the partial
  trace composes, i.e. that $\tr_{\rmaux}[\cdot] = \tr_{\rmaux_B}[\tr_{\rmaux_A}[\cdot]]$.
  \begin{align*}
    &\tr_{\rmaux}[ (M'_a \ot N'_b) (\ket{\psi}\bra{\psi}
    \ot \ket{\rmaux_{A}}\bra{\rmaux_{A}} \ot
      \ket{\rmaux_B}\bra{\rmaux_B}) (M'_a \ot N'_b)]\\
    &\qquad= \tr_{\rmaux_B}[
                                                    \sqrt{M_a \ot I}
                                                    ((I \ot N'_b) \ket{\psi}
                                                    \ket{\rmaux_A}
                                                    \bra{\rmaux_A}
                                                    \bra{\psi} (I \ot
                                                    N'_b)) \sqrt{M_a
                                                    \ot I}] \\
    &\qquad= \sqrt{I \ot N_b} \sqrt{M_a \ot I} \ket{\psi} \bra{\psi}
      \sqrt{M_a \ot I} \sqrt{I \ot N_b}.\qedhere 
  \end{align*}
\end{proof}


\subsection{Pauli matrices and the EPR state}
\label{sec:paulis}
Over a finite field $\F_q$ with order $q = p^t$ for prime $p$, the single-qudit Pauli matrices are a set
of unitary matrices acting on $\mathbb{C}^q$. Every Pauli matrix can
be uniquely written as a product $\omega^{a} X(x) Z(z)$, where $\omega$ is the $p$-th root
$\omega = e^{2\pi/p}$, and $X(x)$ and $Z(z)$ are the matrices
\begin{equation}
  X(x) = \sum_{j \in \F_q} \ket{j+x}\bra{j}, \qquad Z(z) =
  \sum_{j \in \F_q} \omega^{\tr[z j]}\ket{j}
  \bra{j}, \label{eq:pauli-definition} \end{equation}
where the arguments $x, z$ are in $\F_q$, $\tr: \F_q \to \F_p$ is the finite field
trace. The set of all Pauli matrices form a group, known as the Pauli
group or the Weyl-Heisenberg group. For the most part,
in this paper, it will suffice to consider only the group elements of
the form
$\omega^a X(x)$ (``$X$-type'' Paulis) and $Z(z)$ (``$Z$-type''
Paulis). Elements of the form $\omega^a X(x)Z(z)$ for $x, z \neq 0$
are sometimes called ``$Y$-type''.

The eigenvalues of $X(x)$ and $Z(z)$ for
all $x, z$ are powers
of $\omega$, as can be seen from the facts $X(x)^p = Z(z)^p =
I$. Any unitary with this property is known as a \emph{(generalized)
  observable}. Every generalized observable $U$ induces a projective
measurement with $p$ outcomes, corresponding to the $p$ possible eigenvalues of
$U$. As a convenient shorthand, we will refer to performing this projective
measurement as ``measuring $U$.''  In the case of of the $X$ and $Z$
operators, the
eigenvectors $\ket{\tau^X_u}$ and $\ket{\tau^Z_u}$ of $X(1)$ and $Z(1)$ are indexed by elements $u$ of $\F_q$,
with eigenvalue $\tr[u]$; thus, each eigenvalue occurs with
multiplicity $q/p$. Explicitly, they are given by
\begin{equation}\label{eq:pauli-eigenstates}
  \ket{\tau^X_u} = \frac{1}{\sqrt{q}} \sum_{v \in \F_q}
  \omega^{-\tr[uv]} \ket{v},\qquad \ket{\tau^Z_u} = \ket{u}.
  \end{equation}
We denote the projectors onto these eigenvectors by $\tau^X_u$ and
$\tau^Z_u$, respectively. These eigenvectors are also the eigenvectors
of the remaining $X(x), Z(z)$ observables, as shown by the following fact.
\begin{fact}
  For $W \in \{X, Z\}$, the observables $W(v)$ are related to the
  projectors $\tau^W_u$ by
  \begin{align}
    W(v) &= \sum_{u} \omega^{\tr[u \cdot v]}
           \tau^W_u \label{eq:pauli-obs-to-proj} \\
    \tau^W_u &= \E_{\bv} \omega^{-\tr[u \cdot \bv]} W(\bv). \label{eq:pauli-proj-to-obs}
  \end{align}
\end{fact}
\begin{proof}
  We start with~\Cref{eq:pauli-obs-to-proj}. For $W = Z$, the relation
  follows immediately from the definitions. For $W = X$, by
  calculation we have:
  \begin{align*}
    \sum_{u} \omega^{\tr[u \cdot x]} \tau^X_u &= \frac{1}{q} \sum_{u,
                                                v, v'} \omega^{\tr[u
                                                \cdot x]}
                                                \omega^{\tr[u(v -
                                                v')]} \ket{v'}\bra{v}
    \\
    &= \sum_{v, v'} \E_{\bu} \omega^{\tr[\bu \cdot (x + v - v')]}\ket{v'}\bra{v}\\
                                              &= \sum_{v} \ket{v+x}
                                                \bra{v} \\
                                              &= X(x),
  \end{align*}
  where we have applied \Cref{fact:averages-to-zero} in passing from
  the second to the third line.
  Now we show \Cref{eq:pauli-proj-to-obs}:
  \begin{equation*}
    \E_{\bv} \omega^{-\tr[u\cdot \bv]} W(\bv) = \E_{\bv} \sum_{a}
                                         \omega^{-\tr[u \cdot \bv]}
                                         \omega^{\tr[\bv \cdot a]}
                                         \tau^W_{a}
                                         = \E_{\bv} \sum_{a}
                                         \omega^{\tr[(a - u) \cdot \bv]}
                                         \tau^W_{a}
                                       = \tau^W_{u},
  \end{equation*}
  where we first applied \Cref{eq:pauli-obs-to-proj} in the first
  equality, and then used \Cref{fact:averages-to-zero} to perform the
  expectation over $\bv$.
\end{proof}
The Pauli matrices obey the commutation
relation
\begin{equation}\label{eq:commutation-relations}
 X(x)Z(z) = \omega^{-\tr[xz]} Z(z)X(x).
 \end{equation}
This follows directly from \Cref{eq:pauli-definition}. It follows from
this that all of the Pauli matrices (including the $Y$-type matrices)
are generalized observables.

The maximally entangled state, or EPR state, over qudits of dimension
$q$ is the state
\[ \ket{\epr_q} = \frac{1}{\sqrt{q}} \sum_{u \in \F_q} \ket{u} \ot \ket{u}. \]
We will write $\ket{\epr_q^n}$ for $\ket{\epr_q}^{\otimes n}$.
This state obeys the stabilizer relations
\begin{align}
  X(x) \ot X(x) \ket{\epr_q} &= Z(z) \ot Z(-z) \ket{\epr_q} =
                             \ket{\epr_q} \label{eq:epr-stab} \\
  \tau^X_u \ot I \ket{\epr_q} &= I \ot \tau^X_{-u}
                              \ket{\epr_q} \label{eq:epr-stab-x} \\
  \tau^Z_u \ot I \ket{\epr_q} &= I \ot \tau^Z_{u} \ket{\epr_q} \label{eq:epr-stab-z}.
\end{align}
Relations~\eqref{eq:epr-stab-x} and~\eqref{eq:epr-stab-z} imply that
measuring $X(x)$ on both halves of an EPR state will yield two
outcomes $a,b$ satisfying $a = -b$, and measuring $Z(z)$ on both
halves will yield two outcomes $a,b$ that are equal.

In the important special case of finite fields with characteristic $2$
(i.e. $\F_q$ for even $q$), $u = -u$ for all $u \in \F_q$, and thus
measuring any of the $X$ and $Z$ operators on both sides of the state
will always yield the same outcome. 

  
\subsection{State dependent distances}

In this section, we introduce two state-dependent distances.
To motivate them, we first define the consistency game, perhaps the simplest nontrivial two-player game.

\begin{definition}
The \emph{consistency game with question~$x$}, denoted $\game_{\mathrm{con}}(x)$ is defined as follows.
The question~$x$ is given to Alice and Bob, who respond with answers~$\ba$ and~$\ba'$, respectively.
The verifier accepts if~$\ba = \ba'$.
\end{definition}

We will typically play the consistency game when~$\bx$, rather than being a fixed question, is drawn from some distribution.
Our first state-dependent distance quantifies the players` success probability in this case.

\begin{definition}
Let $\{A^x_a\}$ and $\{B^x_a\}$ be sets of matrices in $\mathcal{L}(\mathcal{H}_A)$ and $\mathcal{L}(\mathcal{H}_B)$, respectively.
Let~$\calD$ be a distribution on questions~$x$ and $\ket{\psi}$ be a state in $\mathcal{H}_A \otimes \mathcal{H}_B$.
Consider the game in which the verifier selects $\bx \sim \calD$ and then plays $\game_{\mathrm{con}}(\bx)$.
We say that
\begin{equation*}
A^x_a \otimes I_{\reg{Bob}} \consistency_{\delta} I_{\reg{Alice}} \otimes B^x_a
\end{equation*}
\emph{on state $\ket{\psi}$ and distribution $\calD$} if
Alice and Bob win with probability~$1-O(\delta)$ using the measurements~$A$ and~$B$, respectively.
\end{definition}

We will sometimes leave the state or distribution unspecified, as they are often clear from context.
This distance has a clear operational interpretation.
Our second state-dependent distance,
defined next,
is more analytic.

\begin{definition}
Let $\{Q^x_a\}$ and $\{R^x_a\}$ be sets of matrices in $\mathcal{L}(\mathcal{H})$.
Let $\calD$ be a distribution on the variables~$x$ and $\ket{\psi}$ be a state in $\mathcal{H}$.
Then we say that \emph{$Q^x_a \approx_{\delta} R^x_a$ on state $\ket{\psi}$ and distribution $\calD$} if
\begin{equation*}
\E_{\bx \sim \calD} \sum_a \Vert (Q_{a}^{\bx} - R_{a}^{\bx} ) \ket{\psi}\Vert^2 = O(\delta).
\end{equation*}
\end{definition}

As above, we will sometimes leave the state or distribution unspecified when clear from context.
This is sometimes referred to as \emph{the} state-dependence distance,
whereas our first distance measure is often referred to as the ``consistency".
A typical setting of parameters is $\mathcal{H} = \mathcal{H}_A \otimes \mathcal{H}_B$,
$Q^x_a := A^x_a \otimes I_{\reg{Bob}}$,
and $R^x_a := I_{\reg{Alice}} \otimes B^x_a$.
In this case, we have the following relationship between the two state-dependent distances.

\begin{fact}\label{fact:agreement}
Let $\{A^x_a\}$ and $\{B^x_a\}$ be POVM measurements.
The following two facts hold.
\begin{enumerate}
\item If $A^x_a \otimes I_{\reg{Bob}} \consistency_{\delta} I_{\reg{Alice}} \otimes B^x_a$
	then $A^x_a \otimes I_{\reg{Bob}} \approx_{\delta} I_{\reg{Alice}} \otimes B^x_a$.
\item If $A^x_a \otimes I_{\reg{Bob}} \approx_{\delta} I_{\reg{Alice}} \otimes B^x_a$ \emph{and} $\{A^x_a\}$ and $\{B^x_a\}$ are projective measurements,
	then $A^x_a \otimes I_{\reg{Bob}} \consistency_{\delta} I_{\reg{Alice}} \otimes B^x_a$.\label{item:partial-converse}
\end{enumerate}
\end{fact}
\begin{proof}
Suppose that $A^x_a \otimes I_{\reg{Bob}} \consistency_{\delta} I_{\reg{Alice}} \otimes B^x_a$.
This is equivalent to the statement
\begin{equation}\label{eq:i-derived-a-formula}
\E_{\bx} \sum_a \bra{\psi} A^{\bx}_a \otimes B^{\bx}_a \ket{\psi} \geq 1- O(\delta).
\end{equation}
As a result, using the fact that~$A$ and~$B$ are POVMs,
\begin{align}
\E_{\bx} \sum_a \Vert (A_{a}^{\bx}\otimes I - I \otimes B_{a}^{\bx} ) \ket{\psi}\Vert^2 
& = \E_{\bx} \sum_a \bra{\psi} ((A_{a}^{\bx})^2\otimes I + I \otimes (B_{a}^{\bx})^2 - 2 A_a^{\bx} \otimes B_{a}^{\bx}) \ket{\psi}\nonumber\\
& \leq \E_{\bx} \sum_a \bra{\psi} (A_{a}^{\bx}\otimes I + I \otimes B_{a}^{\bx} - 2 A_a^{\bx} \otimes B_{a}^{\bx}) \ket{\psi}\label{eq:sometimes-tight}\\
& = 2 - 2 \E_{\bx} \sum_a \bra{\psi} A_a^{\bx} \otimes B_{a}^{\bx} \ket{\psi}\nonumber.
\end{align}
By \Cref{eq:i-derived-a-formula}, this is $O(\delta)$.
As a result, $A^x_a \otimes I_{\reg{Bob}} \approx_{\delta} I_{\reg{Alice}} \otimes B^x_a$.
The reverse statement holds when~$A$ and~$B$ are projective measurements
because \Cref{eq:sometimes-tight} is an equality in this case.
\end{proof}

The following fact shows that we can derive a weaker converse in the case when only one of~$A$ or~$B$ is projective.

\begin{fact}\label{fact:almost-agreement}
Suppose $\{A^x_a\}$ and $\{B^x_a\}$ are two measurements such that $A^x_a \otimes I_{\reg{Bob}} \approx_{\delta} I_{\reg{Alice}} \otimes B^x_a$.
Suppose further that either~$A$ or~$B$ is a projective measurement (and the other is a POVM measurement).
Then $A^x_a \otimes I_{\reg{Bob}} \consistency_{\delta^{1/2}} I_{\reg{Alice}} \otimes B^x_a$.
\end{fact}
\begin{proof}
Our goal is to upper bound the expression
\begin{equation}\label{eq:anand-is-a-stable-genius-yet-again}
1 - \E_\bx \sum_a \bra{\psi} A^{\bx}_a \otimes B^{\bx}_a \ket{\psi}.
\end{equation}
We begin by rewriting the number~$1$.
Here we use the fact that because~$A$ is projective, $(A^x_a)^2 = A^x_a$,
and so $\sum_a (A^x_a)^2 = I$.
This gives us:
\begin{align}
\eqref{eq:anand-is-a-stable-genius-yet-again}
& = \E_\bx \sum_a \bra{\psi} (A^{\bx}_a)^2 \otimes I_{\reg{Bob}} \ket{\psi} - \E_\bx \sum_a \bra{\psi} A^{\bx}_a \otimes B^{\bx}_a \ket{\psi}\nonumber\\
& = \E_\bx \sum_a \bra{\psi} (A^{\bx}_a \otimes I_{\reg{Bob}}) \cdot (A^{\bx}_a \otimes I_{\reg{Bob}} - I_{\reg{Alice}} \otimes B^{\bx}_a) \ket{\psi}\nonumber\\
& \leq \E_\bx \sum_a \Vert A^{\bx}_a \otimes I_{\reg{Bob}} \ket{\psi}\Vert \cdot \Vert(A^{\bx}_a \otimes I_{\reg{Bob}} - I_{\reg{Alice}} \otimes B^{\bx}_a) \ket{\psi}\Vert\label{eq:anand-continues-to-be-a-stable-genius}
\end{align}
Now we apply Cauchy-Schwarz and then Jensen's inequality:
\begin{equation*}
\eqref{eq:anand-continues-to-be-a-stable-genius}
\leq  \sqrt{\E_\bx \sum_a \Vert A^{\bx}_a \otimes I_{\reg{Bob}} \ket{\psi}\Vert^2
	\cdot \sum_a \Vert(A^{\bx}_a \otimes I_{\reg{Bob}} - I_{\reg{Alice}} \otimes B^{\bx}_a) \ket{\psi}\Vert^2}
\end{equation*}
The first of these terms we bound by~$1$, and the second is~$O(\delta)$ by assumption.
\end{proof}

\begin{remark}\label{remark:fraction-of-identity}
We note that the requirement in \Cref{fact:almost-agreement}
that one of the two measurements be projective is necessary.
Consider the measurements $\{A^x_a\}$ and $\{B^x_a\}$ with~$m$ separate outcomes~$a$
in which $A^x_a = B^x_a = I/m$ for all~$x$ and~$a$. Then $A^x_a \otimes I_{\reg{Bob}} \approx_{0} I_{\reg{Alice}} \otimes B^x_a$, but as $m \rightarrow \infty$,
\begin{equation*}
1- \E_\bx \sum_a \bra{\psi} A^{\bx}_a \otimes B^{\bx}_a \ket{\psi} \rightarrow 1.
\end{equation*}
\end{remark}

Thus, when one of the measurements is projective,
the ``$\approx_\delta$" distance is roughly equivalent to the ``$\consistency_\delta$" distance,
up to a polynomial factor (which we can tolerate losing in our proofs).
More generally, however, the ``$\approx_\delta$" distance
can be viewed as a weakening of the ``$\consistency_\delta$" distance.
In spite of this, we will spend much of the paper dealing with the ``$\approx_\delta$" distance,
as it is easier to manipulate but still strong enough to reach our desired consequences.
(See \cite[Section~$2.3.1$]{Vid11} for a further defense of this distance.)
We note that even when~$\ket{\psi}$ is a bipartite state,
the ``$\approx_\delta$" distance is defined for matrices~$Q$ and~$R$
which are not necessarily tensor products over the bipartition.
Such matrices will often be useful to pass through during intermediate steps of our proofs.

A common use case for these distances is when the verifier (i) samples a pair of questions~$\bx = (\bx_0, \bx_1)$,
(ii) hands~$\bx_0$ to Alice and~$\bx_1$ second to Bob,
(iii) receives their answers $\ba_0$ and~$\ba_1$
and (iv) accepts if $f(\bx, \ba_0) = g(\bx, \ba_1)$ for some functions~$f$ and~$g$.
Write $\{A^{x_0}_{a_0}\}_{a_0}$ and $\{B^{x_1}_{a_1}\}_{a_1}$ for Alice and Bob's measurements, respectively.
We can view these as measurements which receive the pair~$(\bx_0, \bx_1)$ and simply ignore one coordinate.
Suppose the verifier accepts with probability $1-\delta$.
Using \Cref{not:marginalize}, we can view this as performing the consistency game
between the measurements $A^{x_0}_{[f(x, a_0) = b]}$ and  $B^{x_1}_{[f(x, a_1) = b]}$.
Hence, we can derive the following two facts:
\begin{equation*}
A^{x_0}_{[f(x, a_0) = b]} \otimes I_{\reg{Bob}}
\consistency_{\delta} I_{\reg{Alice}} \otimes B^{x_1}_{[f(x, a_1) = b]},
\quad
\text{and}
\quad
A^{x_0}_{[f(x, a_0) = b]} \otimes I_{\reg{Bob}}
\approx_{\delta} I_{\reg{Alice}} \otimes B^{x_1}_{[f(x, a_1) = b]}.
\end{equation*}
This generic format will be the most common use of these notations.

\Cref{remark:fraction-of-identity} highlights the importance of \emph{projective} measurements when dealing with the~``$\approx_\delta$" distance.
This, and several other key facts about the ``$\approx_\delta$" distance, are true only for projective measurements.
As a result, we will sometimes apply Naimark's theorem (\Cref{thm:naimark})
during our proofs to ``round" POVM measurements into projective measurements.
However, there is a subtlety in doing so, namely that 
because Naimark's theorem preserves measurement outcomes,
any ``$\consistency_\delta$" statements we have derived about our measurement operators will remain true,
but Naimark's theorem is \emph{not} guaranteed to preserve ``$\approx_\delta$" statements.
In this work, we will be able to dispense with this subtlety and assume all ``$\approx_\delta$" statements \emph{are} preserved,
because \emph{all ``$\approx_\delta$" statements in our proofs will be derived from ``$\consistency_\delta$" statements},
and so they will remain true after performing Naimark's theorem, since we could simply rederive them.

\subsection{Miscellaneous properties of the state-dependent distances}

In this section, we record some facts about the ``$\approx_\delta$" notation which we will use repeatedly throughout the paper.
A good rule of thumb is that everything one expects to be true about the ``$\approx_\delta$" notation actually \emph{is}  true,
except for those things which are not. 
As a result, we will be overly pedantic in this section in order to call attention to these cases.

\subsubsection{Simple state-dependent distance facts}

\begin{fact}\label{fact:obvious-vector-fact}
For two vectors $\ket{\psi_1}, \ket{\psi_2}$, 
$\displaystyle
\Vert \ket{\psi_1} + \ket{\psi_2} \Vert^2 \leq 2 \Vert \ket{\psi_1} \Vert^2 + 2 \Vert \ket{\psi_2} \Vert^2.
$
\end{fact}

\begin{fact}\label{fact:measurement-sub-measurement-switcheroo}
Let $\{A_a\}$ be a measurement. Then
\begin{equation*}
\sum_a \Vert A_{a} \ket{\psi}\Vert^2 \leq\Vert \ket{\psi} \Vert^2.
\end{equation*}
\end{fact}
\begin{proof}
If $\{A_a\}$ is a measurement, then
\begin{equation*}
\sum_a \Vert A_{a} \ket{\psi}\Vert^2
= \sum_a \tr( A_a A_{a} \ket{\psi}\bra{\psi})
\leq \tr(I \cdot \ket{\psi}\bra{\psi}) = \Vert \ket{\psi}\Vert^2.\qedhere
\end{equation*}
\end{proof}

\begin{fact}\label{fact:easy-bound}
Let $\{A_a\}$, $\{B_a\}$ be measurements. Then for any state~$\ket{\psi}$,
\begin{equation*}
\sum_a \Vert (A_{a} - B_{a} ) \ket{\psi}\Vert^2 \leq 4 \cdot \Vert \ket{\psi} \Vert^2.
\end{equation*}
\end{fact}
\begin{proof}
By \Cref{fact:obvious-vector-fact},
\begin{equation*}
\sum_a \Vert (A_{a} - B_{a} ) \ket{\psi}\Vert^2 \leq 2 \sum_a \Vert A_{a} \ket{\psi}\Vert^2 + 2 \sum_a \Vert  B_{a} \ket{\psi}\Vert^2.
\end{equation*}
The fact now follows from \Cref{fact:measurement-sub-measurement-switcheroo}.
\end{proof}

\begin{fact}\label{fact:trivial-upper-bound-approx-delta}
Let $\{A^x_a\}$ and $\{B^x_a\}$ be POVM measurements.
Then $A^x_a \approx_1 B^x_a$.
\end{fact}

\begin{fact}\label{fact:add-a-proj}
Let $\{A_a^{x}\}$ and $\{B_a^{x}\}$ be matrices.
Let $\{C_b^{y}\}$ be matrices such that $\sum_b (C_b^y)^\dagger C_b^y \leq I$ for all~$y$.
(This includes the case when $\{C_b^y\}$ form projective or POVM measurements.) Then
\begin{equation*}
\text{$A_a^{x} \approx_{\delta} B_a^x$ implies $C_b^y A_a^x \approx_{\delta} C_b^y B_a^x$.}
\end{equation*}
\end{fact}
\begin{proof}
Fix questions~$x, y$ and answers~$a$. Because of our property on $\{C_b^y\}_b$,
\begin{align*}
\sum_{b} \Vert (C_b^{y} A_a^{x} - C_b^{y} B_a^{x})\ket{\psi}\Vert^2
&= \sum_{b} \bra{\psi} (A_a^x - B_a^x)^\dagger (C_b^y)^\dagger (C_b^y) (A_a^{x} -  B_a^{x})\ket{\psi}\\
&\leq \bra{\psi} (A_a^x - B_a^x)^\dagger (A_a^x - B_a^x) \ket{\psi}
 = \Vert (A_a^{x} -B_a^{x})\ket{\psi}\Vert^2
\end{align*}
We can therefore derive our desired conclusion:
\begin{equation*}
\E_{\bx, \by} \sum_{a, b} \Vert (C_b^{\by} A_a^{\bx} - C_b^{\by} B_a^{\bx}) \ket{\psi} \Vert^2
\leq \E_{\bx, \by} \sum_{a} \Vert (A_a^{\bx} - B_a^{\bx}) \ket{\psi} \Vert^2
= \delta.\qedhere
\end{equation*}
\end{proof}

\begin{fact}\label{fact:swap-dists}
Let $\calD, \calD'$ be two distributions such that $d_{\mathrm{TV}}(\calD, \calD') \leq \epsilon$.
Let $\{A_a^x\}$ and $\{B_a^x\}$ be measurements,
and suppose $A_a^x \approx_\delta B_a^x$ with respect to~$\calD$. Then $A_a^x \approx_{\delta + \epsilon} B_a^x$ with respect to~$\calD'$.
\end{fact}
\begin{proof}
By the definition of total variation distance,
for any set of numbers $\{\nu_x\}$ satisfying $0 \leq \nu_x \leq c$,
the expectations under the two distributions are similar:
\begin{equation*}
|\E_{\bx \sim \calD} [\nu_\bx] - \E_{\bx \sim \calD'}[\nu_{\bx}]| \leq c \cdot \epsilon.
\end{equation*}
We will take for our numbers $\nu_x = \sum_a \Vert (A_{a}^{x} - B_{a}^{x} ) \ket{\psi}\Vert^2$,
which is always less than~$4$ by \Cref{fact:easy-bound}.
As a result,
\begin{equation*}
\E_{\bx \sim \calD'} \nu_{\bx} \leq \E_{\bx \sim \calD} \nu_{\bx} + 4 \eps \leq \delta + 4 \eps.
\end{equation*}
This is $O(\delta + \epsilon)$, which proves the fact.
\end{proof}

\begin{fact}\label{fact:the-ol-state-y-swaperoonie}
Suppose $A_a^x \approx_{\delta} B_a^x$ on state $\ket{\psi}$, and suppose $\Vert \ket{\psi} - \ket{\overline{\psi}} \Vert^2 \leq \epsilon$.
Then $A_a^x \approx_{\delta+\epsilon} B_a^x$ on state $\ket{\overline{\psi}}$.
\end{fact}
\begin{proof}
Applying \Cref{fact:obvious-vector-fact} to $(A_{a}^{x} - B_{a}^{x} ) \ket{\psi}$ and $(A_{a}^{x} - B_{a}^{x} ) (\ket{\overline{\psi}} - \ket{\psi})$,
\begin{equation*}
\E_{\bx \sim \calD} \sum_a \Vert (A_{a}^{\bx} - B_{a}^{\bx} ) \ket{\overline{\psi}}\Vert^2
\leq \E_{\bx \sim \calD} \sum_a \Vert (A_{a}^{\bx} - B_{a}^{\bx} ) \ket{\psi}\Vert^2 + \E_{\bx \sim \calD} \sum_a \Vert (A_{a}^{\bx} - B_{a}^{\bx} ) (\ket{\overline{\psi}}-\ket{\psi})\Vert^2.
\end{equation*}
The first of these is bounded by $O(\delta)$ by the assumption, and the second of these is bounded by $4\epsilon$ by the assumption and \Cref{fact:easy-bound}.
\end{proof}

\begin{fact}\label{fact:close-on-marginal}
Suppose $A_a^{x_1} \approx_{\delta} B_a^{x_1}$ with respect to a distribution $\calD_{\mathrm{margin}}$ on $x_1$.
Let $\calD$ be a distribution on $(x_1, x_2)$ such that the marginal distribution on~$x_1$ is $\calD_{\mathrm{margin}}$.
Then $A_a^{x_1} \approx_{\delta} B_a^{x_1}$ with respect to~$\calD$.
\end{fact}
\begin{proof}
This is a simple calculation involving \Cref{not:marginalize}.
\end{proof}

\begin{fact}\label{fact:average-over-dists}
Let $k$ be a constant, and consider distributions over questions $\calD_1, \ldots, \calD_k$.
Let $\calD$ be a mixture of these distributions,
meaning that there is a probability distribution $p = (p_1, \ldots, p_k)$
such that a draw from $\calD$ can be simulated as follows: draw $\bi \sim p$ and output $\bx$ sampled from $\calD_{\bi}$.
Suppose $A^x_a \approx_{\delta} B^{x}_a$ with respect to $\calD_i$, for all $i \in [k]$.
Then $A^x_a \approx_{\delta} B^{x}_a$ with respect to $\calD$.
\end{fact}
\begin{proof}
By definition, for each $i \in [k]$ there is some constant $C_i$ such that 
\begin{equation*}
\E_{\bx \sim \calD_i} \sum_{a} \Vert (A^{\bx}_a - B^{\bx}_b) \ket{\psi} \Vert^2 \leq C_k \cdot \delta.
\end{equation*}
Then we can bound the mixture with
\begin{equation*}
\E_{\bx \sim \calD} \sum_{a} \Vert (A^{\bx}_a - B^{\bx}_b) \ket{\psi} \Vert^2
= \E_{\bi \sim p} \E_{\bx \sim \calD_{\bi}} \sum_{a} \Vert (A^{\bx}_a - B^{\bx}_b) \ket{\psi} \Vert^2
\leq \E_{\bi \sim p} C_{\bi} \cdot \delta \leq \max_{i}\{C_i\} \cdot \delta = O(\delta).\qedhere
\end{equation*}
\end{proof}

\begin{fact}\label{fact:sub-zero}
Suppose $\{A^x_a\}$ is a projective measurement and $\{B^x_a\}$ is a set of matrices such that
each~$B^x_a$ is positive semidefinite
and $\sum_{a}B^x_a \preceq I$.
Define $C^x_a$ such that for each~$x$, there exists an~$a$ such that $C^x_a := B^x_a + (I - \sum_{a'} B^x_{a'})$
and for all other~$a' \neq a$, $C^x_{a'} := B^x_{a'}$.
Thus, $C^x$ is a POVM for each~$x$.
If $A^x_a \approx_{\eps} B^x_a$ then $A^x_a \approx_{\eps^{1/2}} C^x_a$.
\end{fact}
\begin{proof}
By \Cref{fact:obvious-vector-fact},
\begin{equation*}
\E_{\bx} \sum_a \Vert (A^{\bx}_a - C^{\bx}_a) \ket{\psi} \Vert^2
\leq 2\E_{\bx} \sum_a \Vert (A^{\bx}_a - B^{\bx}_a) \ket{\psi} \Vert^2 + 2 \E_{\bx}\Vert (I - \sum_a B^{\bx}_a) \ket{\psi}\Vert^2.
\end{equation*}
The first of these terms we can bound by $O(\eps)$.  As for the second, 
\begin{multline*}
\E_{\bx}\Vert (I - \sum_a B^{\bx}_a) \ket{\psi}\Vert^2
= \E_{\bx} \bra{\psi} (I-\sum_a B^{\bx}_a)^2 \ket{\psi}
\leq \E_{\bx} \bra{\psi} (I-\sum_a B^{\bx}_a) \ket{\psi}\\
= 1- \E_{\bx} \sum_a\bra{\psi}  B^{\bx}_a \ket{\psi}
\leq1- \E_{\bx} \sum_a\bra{\psi}  (B^{\bx}_a)^2 \ket{\psi}.
\end{multline*}
Now, we write $1 = \E_{\bx} \sum_a \bra{\psi} (A^{\bx}_a)^2 \ket{\psi}$,
which holds because~$A$ is a projective measurement.
We bound the result as follows.
\begin{multline*}
\E_{\bx} \sum_a \bra{\psi} ((A^{\bx}_a)^2  - (B^{\bx}_a)^2)\ket{\psi}
= \Re\left(\E_{\bx} \sum_a \bra{\psi} (A^{\bx}_a  + B^{\bx}_a)(A^{\bx}_a - B^{\bx}_a)\ket{\psi}\right)\\
\leq \E_{\bx} \sqrt{\sum_a \Vert (A^{\bx}_a + B^{\bx}_a) \ket{\psi}\Vert^2}
	\cdot \sqrt{\sum_a \Vert (A^{\bx}_a - B^{\bx}_a)\ket{\psi} \Vert^2}.
\end{multline*}
For each~$\bx$,
we can bound the first square root by $O(1)$
due to \Cref{fact:obvious-vector-fact} and \Cref{fact:measurement-sub-measurement-switcheroo}.
Having done so, we can move the expectation into the second square root by Jensen's inequality.
The result is $O(\eps^{1/2})$ by assumption. This proves the fact.
\end{proof}

\subsubsection{Data processing}

In this section, we show a simple data processing inequality for the ``$\consistency_\delta$" distance.
We also observe that one does \emph{not} hold for the ``$\approx_\delta$" distance.

\begin{fact}\label{fact:specialize-the-simeq}
Suppose that $A^x_a \otimes I_{\reg{Bob}} \consistency_\delta I_{\reg{Alice}} \otimes B^x_a$.
Then $A^x_{[f(a) = b]} \otimes I_{\reg{Bob}} \consistency_\delta I_{\reg{Alice}} \otimes B^x_{[f(a) = b]}$.
\end{fact}
\begin{proof}
Given question~$\bx$, if Alice and Bob return~$\ba$ and~$\ba'$ in which $\ba = \ba'$, then $f(\ba) = f(\ba')$.
As a result, applying~$f$ to their answers cannot decrease the probability they agree.
\end{proof}

\begin{remark}
We note that the same fact is \emph{not} true for the ``$\approx_\delta$" distance.
Consider answers of the form $a = (b, i)$, where $b \in \{0, 1\}$ and $i \in [m]$.
Suppose $A^x_{b, i} = I/(2m)$ for all~$a$, whereas $B^x_{0, i} = I/m$ and $B^x_{1, i} = 0$ for all~$i$.
Consider the function $f(b, i) = b$.
It can be checked that in this case, $A^x_a \otimes I_{\reg{Bob}} \approx_{1/2m} I_{\reg{Alice}} \otimes B^x_a$
but $A^x_{[f(a) = b]} \otimes I_{\reg{Bob}} \approx_{1/2} I_{\reg{Alice}} \otimes B^x_{[f(a) = b]}$.
\end{remark}

\subsubsection{Triangle inequalities}

In this section, we give two triangle inequalities. Our first is for the state-dependent distance.

\begin{fact}[Triangle inequality]\label{fact:triangle}
Suppose $A_a^x \approx_\delta B_a^x$ and $B_a^x \approx_\epsilon C_a^x$. Then $A_a^x \approx_{\delta + \epsilon} C_a^x$.
\end{fact}
\begin{proof}
Applying \Cref{fact:obvious-vector-fact} to $(A_{a}^{x} - B_{a}^{x} ) \ket{\psi}$ and $(B_{a}^{x} - C_{a}^{x} ) \ket{\psi}$,
\begin{align*}
\E_{\bx \sim \calD} \sum_a \Vert (A_{a}^{\bx} - C_{a}^{\bx} ) \ket{\psi}\Vert^2
&\leq 2 \E_{\bx \sim \calD} \sum_a \Vert (A_{a}^{\bx} - B_{a}^{\bx} ) \ket{\psi}\Vert^2 + 2 \E_{\bx \sim \calD} \sum_a \Vert (B_{a}^{\bx} - C_{a}^{\bx} ) \ket{\psi}\Vert^2\\
&\leq 2 (\delta + \epsilon).\qedhere
\end{align*}
\end{proof}

Note that this does \emph{not} show that if
\begin{equation*}
A^x_a \otimes I_{\reg{Bob}} \approx_\delta I_{\reg{Alice}} \otimes B^x_a
\quad\text{and}\quad
B^x_a \otimes I_{\reg{Bob}} \approx_\delta I_{\reg{Alice}} \otimes C^x_a
\end{equation*}
then $A^x_a \otimes I_{\reg{Bob}} \approx_\delta I_{\reg{Alice}} \otimes C^x_a$.
This would only follow if, for example, we also knew that $D^x_a \otimes I_{\reg{Bob}} \approx_\delta I_{\reg{Alice}} \otimes D^x_a$,
for $D$ equal to one of~$A$, $B$, or $C$.
We do, however, always have the following triangle-like inequalities.

\begin{fact}[Triangle-like inequalities]\label{fact:triangle-like}
The following two facts are true.
\begin{enumerate}
\item Suppose $A^x_a \otimes I_{\reg{Bob}} \consistency_\delta I_{\reg{Alice}} \otimes B^x_a$,
	$B^x_a \otimes I_{\reg{Bob}} \consistency_\delta I_{\reg{Alice}} \otimes C^x_a$,
	and $C^x_a \otimes I_{\reg{Bob}} \consistency_\delta I_{\reg{Alice}} \otimes D^x_a$.
	Then $A^x_a \otimes I_{\reg{Bob}} \consistency_\delta I_{\reg{Alice}} \otimes D^x_a$.
\item Suppose $A^x_a \otimes I_{\reg{Bob}} \approx_\delta I_{\reg{Alice}} \otimes B^x_a$,
	$B^x_a \otimes I_{\reg{Bob}} \approx_\delta I_{\reg{Alice}} \otimes C^x_a$,
	and $C^x_a \otimes I_{\reg{Bob}} \approx_\delta I_{\reg{Alice}} \otimes D^x_a$.
	Then $A^x_a \otimes I_{\reg{Bob}} \approx_\delta I_{\reg{Alice}} \otimes D^x_a$.
\end{enumerate}
\end{fact}

Before proving this, we need the following fact from linear algebra.

\begin{fact}\label{fact:fun-matrix-fact}
Suppose $0 \preceq A, B, C, D \preceq I$.  Then
\begin{equation*}
1- \bra{\psi} A \otimes D \ket{\psi}
\leq (1 - \bra{\psi} A \otimes B \ket{\psi}) + (1 - \bra{\psi} B \otimes C \ket{\psi}) + (1 - \bra{\psi} C \otimes D \ket{\psi}).
\end{equation*}
\end{fact}
\begin{proof}
Rearranging, we want to show that
\begin{equation*}
\bra{\psi} A \otimes B \ket{\psi} + \bra{\psi} B \otimes C \ket{\psi} + \bra{\psi} C \otimes D \ket{\psi} - \bra{\psi} A \otimes D \ket{\psi}
\leq 2.
\end{equation*}
Or, equivalently
\begin{equation*}
\tr(\ket{\psi}\bra{\psi} \cdot (A \otimes B + B \otimes C + C \otimes D - A \otimes D)) \leq 2.
\end{equation*}
The left-hand side is at most the maximum eigenvalue of $A \otimes B + B \otimes C + C \otimes D - A \otimes D$.
To bound this maximum eigenvalue, we note that $A \otimes B \preceq A \otimes I$, $B \otimes C \preceq I \otimes I$, and $C \otimes D \preceq I \otimes D$.
As a result, 
\begin{equation*}
A \otimes B + B \otimes C + C \otimes D - A \otimes D
\preceq A \otimes I + I \otimes I + I \otimes D - A \otimes D.
\end{equation*}
Next, $I \otimes D - A \otimes D = (I - A) \otimes D \preceq (I-A) \otimes I$ because $A \preceq I$.
Thus,
\begin{equation*}
A \otimes I + I \otimes I + I \otimes D - A \otimes D
\preceq A \otimes I + I \otimes I  + (I-A) \otimes I = 2 \cdot I \otimes I.
\end{equation*}
But the maximum eigenvalue of this is~$2$.
\end{proof}


Now we prove \Cref{fact:triangle-like}.
\begin{proof}[Proof of \Cref{fact:triangle-like}]
The second fact follows from several applications of \Cref{fact:triangle}.
As for the first fact, we can write the consistency as
\begin{equation*}
\E_\bx \sum_a ( 1- \bra{\psi} A^{\bx}_a \otimes D^{\bx}_a \ket{\psi})
\end{equation*}
Applying \Cref{fact:fun-matrix-fact}, this is at most
\begin{equation*}
\E_\bx \sum_a (1 - \bra{\psi} A^\bx_a \otimes B^\bx_a \ket{\psi}) + (1 - \bra{\psi} B^\bx_a \otimes C^\bx_a \ket{\psi}) + (1 - \bra{\psi} C^\bx_a \otimes D^\bx_a \ket{\psi})
\end{equation*}
Averaging over questions and summing over answers, each of these terms is at most~$\delta$, by assumption.
\end{proof}

\subsubsection{Close strategies have close game values}

In this section, we will show that two strategies which are close in state-dependent distance are also close in value for any game~$\game$.
We note crucially that one of the two strategies must be \emph{projective} to apply this fact.

\begin{fact}\label{fact:approx-delta-generalized-game-value}
Let $\calD$ be a distribution on questions~$x$, and for each~$x$ let $\mathrm{acc}(x)$ be a set of ``accepting" answers.
Given a state~$\psi$ and a strategy $\{A^x_a\}$ define
\begin{equation*}
\mathrm{val}(A) = \E_{\bx \sim \calD} \sum_{a \in \mathrm{acc}(x)} \bra{\psi} A^{\bx}_a \ket{\psi}.
\end{equation*}
Suppose $\{A^x_a\}$ and $\{B^x_a\}$ are two strategies such that $A^x_a \approx_\delta B^x_a$ on state~$\psi$ and distribution~$\calD$.
Suppose further that either~$A$ or~$B$ is a projective measurement (and the other is a POVM measurement).
Then
\begin{equation*}
\mathrm{val}(A) - O(\delta^{1/2}) \leq \mathrm{val}(B) \leq \mathrm{val}(A) + O(\delta^{1/2}).
\end{equation*}
\end{fact}
\begin{proof}
Assume without loss of generality that~$A$ is a projective measurement and~$B$ is a POVM measurement.
We will prove the fact by showing the following stronger statement:
for each~$x$, let~$S(x)$ be any set of answers~$a$,
and define
\begin{equation*}
\mathrm{val}(A,S) := \E_\bx \sum_{a \in S(x)} \bra{\psi} A^{\bx}_a\ket{\psi}.
\end{equation*}
Then $\mathrm{val}(A,S) \leq \mathrm{val}(B,S) + O(\delta^{1/2})$.
By taking $S(x) := \mathrm{acc}(x)$ this implies the lower bound $\mathrm{val}(A) - O(\delta^{1/2}) \leq \mathrm{val}(B)$,
and by taking $S(x) := \mathrm{rej}(x)$, defined to be the set of answers \emph{not} in $\mathrm{acc}(x)$,
then this implies the upper bound $\mathrm{val}(B) \leq \mathrm{val}(A) + O(\delta^{1/2})$.

If we write $\ket{u_a^x} = A_{a}^{x}  \ket{\psi}$
	and $\ket{w_a^x} = (B_{a}^{x} - A_{a}^{x}) \ket{\psi}$, then
\begin{equation*}
\Vert B_{a}^{x} \ket{\psi} \Vert^2
= \Vert \ket{u_a^x} + \ket{v_a^x} \Vert^2
 = \Vert\ket{u_a^x}\Vert^2 + \Vert\ket{w_a^x}\Vert^2
 + \braket{u_a^x \mid w_a^x} + \braket{w_a^x \mid u_a^x}.
\end{equation*}
By definition,
\begin{align*}
\mathrm{val}(B)
& = \E_{\bx} \sum_{a\in S(\bx)} \bra{\psi} B_{a}^{\bx} \ket{\psi}\\
& \geq \E_{\bx} \sum_{a\in S(\bx)} \bra{\psi} (B_{a}^{\bx})^2 \ket{\psi}\tag{because~$B$ is a POVM}\\
& = \E_{\bx} \sum_{a\in S(\bx)}\Vert B_{a}^{\bx} \ket{\psi} \Vert^2\\
& = \E_{\bx} \sum_{a\in S(\bx)} \Vert\ket{u_a^{\bx}}\Vert^2 + \Vert\ket{w_a^{\bx}}\Vert^2
		 + \braket{u_a^{\bx} \mid w_a^{\bx}} + \braket{w_a^{\bx} \mid u_a^{\bx}}.
\end{align*}
Averaging over questions and summing over answers, the first term is exactly $\mathrm{val}(A)$ because~$A$ is projective.
The second term is always nonnegative, so we lower bound it by zero.
As for the last two terms,
\begin{equation}\label{eq:our-powers-combined--wind!}
\braket{u_a^x \mid w_a^x} + \braket{w_a^x \mid u_a^x}
\geq - 2 \cdot |\braket{u_a^x \mid w_a^x}| 
\geq -2 \cdot \Vert u_a^x \Vert \cdot \Vert w_a^x \Vert.
\end{equation}
Applying Cauchy-Schwarz,
Jensen's inequality,
and \Cref{fact:measurement-sub-measurement-switcheroo},
\begin{equation}\label{eq:our-powers-combined--heart!}
\E_{\bx} \sum_{a \in S(\bx)} \Vert u_a^{\bx} \Vert \cdot \Vert w_a^{\bx} \Vert
\leq \E_{\bx} \sqrt{\sum_{a \in S(\bx)} \Vert u_a^{\bx} \Vert^2 \cdot  \sum_{a \in S(\bx)}\Vert w_a^{\bx}\Vert^2 }
\leq \sqrt{\E_{\bx} \sum_{a \in S(\bx)}\Vert w_a^{\bx}\Vert^2 }.
\end{equation}
But the expectation inside the root is at most $O(\delta)$ because $A^{x}_a \approx_{\delta} B^{x}_a$.
Combining \Cref{eq:our-powers-combined--wind!,eq:our-powers-combined--heart!} completes the proof.
\end{proof}

We will typically, though not always,
apply \Cref{fact:approx-delta-generalized-game-value} in the following special case.

\begin{fact}\label{fact:approx-delta-game-value}
Let $\game$ be a game whose questions $(\bx_1, \bx_2) \sim \game$ have marginal distribution $\bx_1 \sim \calD$.
Suppose $\{A_a^x\}$
and $\{B_a^x\}$ are measurements such that
$
A_a^x  \otimes I \approx_\delta B_a^x \otimes I
$
on state $\psi$ and distribution~$\calD$.
Consider the strategies $\calS_A = \{\psi, A\}$ and $\calS_B = \{\psi, B\}$.
If either~$A$ or~$B$ is a projective measurement (and the other is a POVM measurement), then
\begin{equation*}
\valstrat{\game}{\calS_A} - O(\delta^{1/2}) \leq \valstrat{\game}{\calS_B} \leq \valstrat{\game}{\calS_A} + O(\delta^{1/2}).
\end{equation*}
\end{fact}
\begin{proof}
First, we observe that
\begin{equation*}
A^{x_1}_{a_1} \otimes A^{x_2}_{a_2}
\approx_{\delta} A^{x_1}_{a_1} \otimes B^{x_2}_{a_2}
\approx_{\delta} B^{x_1}_{a_1} \otimes B^{x_2}_{a_2}
\end{equation*}
by \Cref{fact:add-a-proj}.
The result follows by applying \Cref{fact:approx-delta-game-value} with
``$A$" set to $A^{x_1}_{a_1} \otimes A^{x_2}_{a_2}$,
``$B$" set to $B^{x_1}_{a_1} \otimes B^{x_2}_{a_2}$,
and ``$\calD$" set to the distribution on $(\bx_1, \bx_2)$.
We note that ``$\mathrm{val}(A)$" there is equal to $\valstrat{\game}{\calS_A}$ here and
``$\mathrm{val}(B)$" there is equal to $\valstrat{\game}{\calS_B}$ here.
\end{proof}

\subsubsection{Generating new measurements}

In this section, we show how to combine multiple measurements into a single measurement by ``sandwiching" them together.

\begin{fact}\label{fact:sandwich}
Let $k \geq 0$ be a constant.
Let $\{A^x_{a_1, \ldots, a_k}\}$ be a projective measurement.
For each $1 \leq i \leq k$, let $\{(B_i)^x_{a_i}\}$ be a projective measurement, and suppose that
\begin{equation}\label{eq:no-take-only-throw}
(A^x_{a_i})_{\reg{Alice}} \otimes I_{\reg{Bob}} \simeq_{\delta} I_{\reg{Alice}} \otimes ((B_i)^x_{a_i})_{\reg{Bob}}.
\end{equation}
Define the POVM measurement $\{J^x_{g_1, \ldots, g_k}\}$ as
\begin{equation*}
J^x_{a_1, \ldots, a_k} := (B_k)^x_{a_k} \cdots (B_2)^x_{a_2} \cdot (B_1)^x_{a_1} \cdot (B_2)^x_{a_2} \cdots (B_k)^x_{a_k}.
\end{equation*}
Then
\begin{equation*}
(A^x_{a_1, \ldots, a_k})_{\reg{Alice}} \otimes I_{\reg{Bob}} \simeq_{\delta^{1/2}} I_{\reg{Alice}} \otimes (J^x_{a_1, \ldots, a_k})_{\reg{Bob}}.
\end{equation*}
\end{fact}
\begin{proof}
For each $1 \leq i \leq k$, 
\Cref{eq:no-take-only-throw} implies that
\begin{equation*}
(A^x_{a_i})_{\reg{Alice}} \otimes I_{\reg{Bob}} \approx_{\delta} I_{\reg{Alice}} \otimes ((B_i)^x_{a_i})_{\reg{Bob}}.
\end{equation*}
Now, we repeatedly apply this using \Cref{fact:add-a-proj}:
\begin{align*}
(A^x_{a_1, \ldots, a_k})_{\reg{Alice}} \otimes I_{\reg{Bob}}
& = (A^x_{a_k} \cdots A^x_{a_2} \cdot A^x_{a_1} \cdot A^x_{a_2} \cdots A^x_{a_{k}})_{\reg{Alice}} \otimes I_{\reg{Bob}}\\
& \approx_{\delta} (A^x_{a_k} \cdots A^x_{a_2} \cdot A^x_{a_1} \cdot A^x_{a_2} \cdots A^x_{a_{k-1}})_{\reg{Alice}} \otimes ((B_k)^x_{a_k})_{\reg{Bob}}\\
&\qquad\qquad\qquad \qquad \qquad \qquad \cdots\\
& \approx_{\delta} I_{\reg{Alice}} \otimes ((B_k)^x_{a_k} \cdots (B_2)^x_{a_2} \cdot (B_1)^x_{a_1} \cdot (B_2)^x_{a_2} \cdots (B_k)^x_{a_k})_{\reg{Bob}}\\
& = I_{\reg{Alice}} \otimes (J^x_{a_1, \ldots, a_k})_{\reg{Bob}}.
\end{align*}
The fact now follows from \Cref{fact:almost-agreement} and the fact that~$A$ is a projective measurement.
\end{proof}

Next, we extend \Cref{fact:sandwich} to the case of polynomial measurements (see \Cref{sec:classical-q-low-deg} below).
These are structured measurements in which the prover returns the evaluation of a function sampled independently from their input.
The goal is to retain this structure even after ``sandwiching" them together.

\begin{fact}\label{fact:low-degree-sandwich}
Let $k \geq 0$ be a constant.
Let $\calD$ be a distribution on questions $x \in \calX$.
For each $1 \leq i \leq k$, let $\calG_i$ be a set of functions $g_i : \calX \rightarrow \calR_i$.
and let $\{G^i_g\}$ be a projective measurement with outcomes from this set.
Suppose that the set $\calG_i$ has the following distance property:
for any two nonequal $g_i, g_i' \in \calG_i$, the probability that $g_i(\bx) = g_i'(\bx)$, over a random $\bx \sim \calD$, is at most $\eps$.

Let $\{A_{g_1, \ldots, g_k}\}$ be a projective measurement with outcomes $g_i \in \calF_i$.
For each $1 \leq i \leq k$, suppose that
\begin{equation}\label{eq:treat-yo-self}
(A_{[g_i(x) = a_i]})_{\reg{Alice}} \otimes I_{\reg{Bob}} \simeq_{\delta} I_{\reg{Alice}} \otimes (G^i_{[g_i(x)=a_i]})_{\reg{Bob}}.
\end{equation}
Define the POVM measurement $\{J_{g_1, \ldots, g_k}\}$ as
\begin{equation*}
J_{g_1, \ldots, g_k} := G^k_{g_k} \cdots G^2_{g_2} \cdot G^1_{g_1} \cdot G^2_{g_2} \cdots G^k_{g_k}.
\end{equation*}
Then
\begin{equation*}
(A_{[g_1(x), \ldots, g_k(x) = a_1, \ldots, a_k]})_{\reg{Alice}} \otimes I_{\reg{Bob}} \simeq_{(\delta+\eps)^{1/2}} I_{\reg{Alice}} \otimes (J_{[g_1(x), \ldots, g_k(x) = a_1, \ldots, a_k]})_{\reg{Bob}}.
\end{equation*}
\end{fact}
\begin{proof}
Let $1 \leq i \leq k$.
By~\Cref{eq:treat-yo-self},
if Alice measures with~$A$, producing~$\bg_i$, and Bob measures with~$G^i$, producing~$\bg_i'$,
then the probability that $\bg_i(\bx) \neq \bg_i'(\bx)$ is $O(\delta)$.
Write $\eta$ for the probability that $\bg_i \neq \bg_i'$.
Then we have the expression $\eta \cdot (1-\eps) \leq O(\delta)$ or, equivalently, $\eta \leq O(\delta/(1-\eps))$.
When $\eps < 1/2$, this gives the bound $\eta \leq O(\delta)$, and when $\eps \geq 1/2$, we have the trivial bound $\eta \leq O(\eps)$.
As a result, $\eta = O(\delta + \eps)$.

In conclusion,
\begin{equation*}
(A_{g_i})_{\reg{Alice}} \otimes I_{\reg{Bob}} \simeq_{\delta + \eps} I_{\reg{Alice}} \otimes (G^i_{g_i})_{\reg{Bob}}.
\end{equation*}
We can now apply \Cref{fact:sandwich} to $A_{g_1, \ldots, g_k}$ and the $G^i_{g_i}$ measurements.
It implies that
\begin{equation*}
(A_{g_1, \ldots, g_k})_{\reg{Alice}} \otimes I_{\reg{Bob}} \simeq_{(\delta+\eps)^{1/2}} I_{\reg{Alice}} \otimes (J_{g_1, \ldots, g_k})_{\reg{Bob}}.
\end{equation*}
The fact now follows from the data processing inequality \Cref{fact:specialize-the-simeq}.
\end{proof}

In our next fact, we show that \Cref{fact:low-degree-sandwich} holds even when we drop the structured assumption on the~$A$ matrix.
The tradeoff is that we must now assume that the~$k$ different measurements act on different parts of the input string.
In this case, the distance condition becomes slightly more cumbersome to state.

\begin{fact}\label{fact:low-degree-sandwich-on-steroids}
Let $k \geq 0$ be a constant.
Let $\calD$ be a distribution on questions $(x, y_1, \ldots, y_k)$, where each $y_i \in \calY_i$.
For each $1 \leq i \leq k$, let $\calG_i$ be a set of functions $g_i : \calY_i \rightarrow \calR_i$.
and let $\{(G_i)^x_g\}$ be a projective measurement with outcomes from this set.
(For the $i=1$ case, we also allow this measurement to be a POVM.)
Suppose that the set $\calG_i$ has the following distance property:
fix a question $z = (x, y_1, \ldots, y_{i-1}, y_{i+1}, \ldots, y_k)$, and let $\calD_z$ be the distribution on~$y_i$ conditioned on the other outcomes~$z$.
Then for any two nonequal $g_i, g_i' \in \calG_i$, the probability that $g_i(\by_i) = g_i'(\by_i)$, over a random $\by_i \sim \calD_z$, is at most $\eps$.

Let $\{A^{x, y_1, \ldots, y_k}_{a_1, \ldots, a_k} \}$ be a projective measurement with outcomes $g_i \in \calF_i$.
For each $1 \leq i \leq k$, suppose that
\begin{equation}\label{eq:gonna-compare-like-a-pear}
(A^{x, y_1, \ldots, y_k}_{a_i})_{\reg{Alice}} \otimes I_{\reg{Bob}} \simeq_{\delta} I_{\reg{Alice}} \otimes ((G_i)^x_{[g_i(y_i)=a_i]})_{\reg{Bob}}.
\end{equation}
Suppose also that
\begin{equation}\label{eq:spend-some-time-reflecting}
(A^{x, y_1, \ldots, y_k}_{a_1, \ldots, a_k})_{\reg{Alice}} \otimes I_{\reg{Bob}}
	\simeq_{\delta} I_{\reg{Alice}} \otimes (A^{x, y_1, \ldots, y_k}_{a_1, \ldots, a_k})_{\reg{Bob}}.
\end{equation}
Define the POVM measurement $\{J^x_{g_1, \ldots, g_k}\}$ as
\begin{equation*}
J^x_{g_1, \ldots, g_k} := (G_k)^x_{g_k} \cdots (G_2)^x_{g_2} \cdot (G_1)^x_{g_1} \cdot (G_2)^x_{g_2} \cdots (G_k)^x_{g_k}.
\end{equation*}
Then
\begin{equation*}
(A^{x, y_1, \ldots, y_k}_{a_1, \ldots, a_k})_{\reg{Alice}} \otimes I_{\reg{Bob}}
	\simeq_{\poly(\delta, \eps)} I_{\reg{Alice}} \otimes (J^x_{[g_1(y_1), \ldots, g_k(y_k) = a_1, \ldots, a_k]})_{\reg{Bob}}.
\end{equation*}
\end{fact}
\begin{proof}
First, we show how to reduce this to the $k = 2$ case. Then we prove it for that case.
Assume the fact holds when $k = 2$.  Define the POVM measurement $\{(J_i)^x_{g_1, \ldots, g_i}\}$ as
\begin{equation*}
(J_i)^x_{g_1, \ldots, g_i} := (G_i)^x_{g_i} \cdots (G_2)^x_{g_2} \cdot (G_1)^x_{g_1} \cdot (G_2)^x_{g_2} \cdots (G_i)^x_{g_i}.
\end{equation*}
We will show by induction that 
\begin{equation}\label{eq:hi-i-work-in-graphics-but-not-that-graphics}
(A^{x, y_1, \ldots, y_k}_{a_1, \ldots, a_i})_{\reg{Alice}} \otimes I_{\reg{Bob}}
	\simeq_{\poly(\delta, \eps)} I_{\reg{Alice}} \otimes ((J_i)^x_{[g_1(y_1), \ldots, g_i(y_i) = a_1, \ldots, a_i]})_{\reg{Bob}},
\end{equation}
the base case being trivial. Assume this holds for~$i$.
We apply the $k = 2$ case as follows:
consider the question tuple $(y_1, \ldots, y_i)$ as a single question
and consider functions of the form $(y_1, \ldots, y_i) \mapsto (g_1(y_1), \ldots, g_i(y_i))$.
Then the first POVM measurement is $J_i$, which satisfies \Cref{eq:gonna-compare-like-a-pear} due to \Cref{eq:hi-i-work-in-graphics-but-not-that-graphics}.
The second measurement is the projector $G_{i+1}$.
Then the $k = 2$ case immediately implies the $i+1$ case of \Cref{eq:hi-i-work-in-graphics-but-not-that-graphics}.

Now we prove the $k = 2$ case.
Our goal is to show that
\begin{equation}\label{eq:food-rakes}
\E_{\bx, \by_1, \by_2} \sum_{a_1, g_2} \bra{\psi} (A^{\bx, \by_1, \by_2}_{a_1, g_2(\by_2)})_{\reg{Alice}}
\otimes ((G_2)^{\bx}_{g_2} \cdot (G_1)^{\bx}_{[g_1(\by_1) = a_1]} \cdot (G_2)^{\bx}_{g_2})_{\reg{Bob}} \ket{\psi}
\end{equation}
is at least $1- \poly(\delta, \eps)$.
We will do this by showing that
\begin{equation}\label{eq:what-a-johnnie-wants}
((G_1)^x_{[g_1(y_1)=a_1]} \cdot (G_2)^x_{g_2})_{\reg{Alice}} \otimes I_{\reg{Bob}}
\approx_{\poly(\delta, \eps)} ((G_2)^x_{g_2} \cdot (G_1)^x_{[g_1(y_1)=a_1]})_{\reg{Alice}} \otimes I_{\reg{Bob}}.
\end{equation}
is at most $\poly(\delta, \eps)$.
To see that this is sufficient, note that the related expression
\begin{equation*}
\E_{\bx, \by_1, \by_2} \sum_{a_1, g_2} \bra{\psi} (A^{\bx, \by_1, \by_2}_{a_1, g_2(\by_2)})_{\reg{Alice}}
\otimes ((G_2)^{\bx}_{g_2} \cdot (G_2)^{\bx}_{g_2} \cdot (G_1)^{\bx}_{[g_1(\by_1) = a_1]})_{\reg{Bob}} \ket{\psi}
\end{equation*}
is exactly equal to~$1$ because $G_2$ is a projector.
Taking the difference between this and \Cref{eq:food-rakes}, we get
\begin{equation*}
\E_{\bx, \by_1, \by_2} \sum_{a_1, g_2} \bra{\psi} (A^{\bx, \by_1, \by_2}_{a_1, g_2(\by_2)})_{\reg{Alice}}
\otimes ((G_2)^{\bx}_{g_2} \cdot ((G_1)^{\bx}_{[g_1(\by_1) = a_1]} \cdot (G_2)^{\bx}_{g_2}- (G_2)^{\bx}_{g_2} \cdot (G_1)^{\bx}_{[g_1(\by_1) = a_1]}))_{\reg{Bob}} \ket{\psi}.
\end{equation*}
Cauchy-Schwarz allows us to bound this by
\begin{multline*}
\leq \E_{\bx, \by_1, \by_2} \sqrt{ \sum_{a_1, g_2}\Vert (A^{\bx, \by_1, \by_2}_{a_1, g_2(\by_2)})_{\reg{Alice}} \otimes ((G_2)^{\bx}_{g_2})_{\reg{Bob}} \ket{\psi} \Vert^2}\\
	\cdot \sqrt{ \sum_{a_1, g_2}\Vert I_{\reg{Alice}} \otimes ((G_1)^{\bx}_{[g_1(\by_1) = a_1]} \cdot (G_2)^{\bx}_{g_2}- (G_2)^{\bx}_{g_2} \cdot (G_1)^{\bx}_{[g_1(\by_1) = a_1]})_{\reg{Bob}} \ket{\psi} \Vert^2}.
\end{multline*}
The expression inside the first square root is always at most~$1$.
This allows us to bring the expectation into the second square root by Jensen's inequality,
and the resulting expression we can bound due to \Cref{eq:what-a-johnnie-wants}.

Now we bound \Cref{eq:what-a-johnnie-wants}.
Showing this is small is equivalent to showing
\begin{equation*}
\E_{\bx, \by_1, \by_2}
\sum_{a_1, g_2}\Vert ((G_1)^{\bx}_{[g_1(\by_1) = a_1]}
	\cdot (G_2)^{\bx}_{g_2}- (G_2)^{\bx}_{g_2} \cdot (G_1)^{\bx}_{[g_1(\by_1) = a_1]})_{\reg{Alice}} \otimes I_{\reg{Bob}} \ket{\psi} \Vert^2
\end{equation*}
is small.  Expanding this, we get
\begin{align}
\E_{\bx, \by_1, \by_2} \sum_{a_1, g_2} \bra{\psi}\big(
	&(G_2)^{\bx}_{g_2} \cdot (G_1)^{\bx}_{[g_1(\by_1) = a_1]} \cdot (G_1)^{\bx}_{[g_1(\by_1) = a_1]} \cdot (G_2)^{\bx}_{g_2} \otimes I_{\reg{Bob}}\nonumber\\
	+ &(G_1)^{\bx}_{[g_1(\by_1) = a_1]} \cdot (G_2)^{\bx}_{g_2} \cdot (G_2)^{\bx}_{g_2} \cdot (G_1)^{\bx}_{[g_1(\by_1) = a_1]} \otimes I_{\reg{Bob}}\nonumber\\
	- &(G_2)^{\bx}_{g_2} \cdot (G_1)^{\bx}_{[g_1(\by_1) = a_1]} \cdot (G_2)^{\bx}_{g_2} \cdot (G_1)^{\bx}_{[g_1(\by_1) = a_1]} \otimes I_{\reg{Bob}}\nonumber\\
	- &(G_1)^{\bx}_{[g_1(\by_1) = a_1]} \cdot (G_2)^{\bx}_{g_2} \cdot (G_1)^{\bx}_{[g_1(\by_1) = a_1]} \cdot (G_2)^{\bx}_{g_2} \otimes I_{\reg{Bob}} \big)\ket{\psi}.\label{eq:burger-king}
\end{align}
We \emph{do} know that $G_1$ and $G_2$ satisfy some form of commutation.
Because they satisfy \Cref{eq:gonna-compare-like-a-pear} and~$A$ is a projector, 
we know that
\begin{equation*}
((G_1)^x_{[g_1(y_1)=a_1]} \cdot (G_2)^x_{[g_2(y_2) = a_2]})_{\reg{Alice}} \otimes I_{\reg{Bob}}
\approx_{\delta} ((G_2)^x_{[g_2(y_2) = a_2]} \cdot (G_1)^x_{[g_1(y_1)=a_1]})_{\reg{Alice}} \otimes I_{\reg{Bob}}.
\end{equation*}
Expanding this as above, we can bound the following expression by~$\delta$:
\begin{align}
\E_{\bx, \by_1, \by_2} \sum_{a_1, a_2} \bra{\psi}\big(
	&(G_2)^\bx_{[g_2(\by_2) = a_2]} \cdot (G_1)^{\bx}_{[g_1(\by_1) = a_1]} \cdot (G_1)^{\bx}_{[g_1(\by_1) = a_1]} \cdot (G_2)^\bx_{[g_2(\by_2) = a_2]} \otimes I_{\reg{Bob}}\nonumber\\
	+ &(G_1)^{\bx}_{[g_1(\by_1) = a_1]} \cdot (G_2)^\bx_{[g_2(\by_2) = a_2]} \cdot (G_2)^\bx_{[g_2(\by_2) = a_2]} \cdot (G_1)^{\bx}_{[g_1(\by_1) = a_1]} \otimes I_{\reg{Bob}}\nonumber\\
	- &(G_2)^\bx_{[g_2(\by_2) = a_2]} \cdot (G_1)^{\bx}_{[g_1(\by_1) = a_1]} \cdot (G_2)^\bx_{[g_2(\by_2) = a_2]} \cdot (G_1)^{\bx}_{[g_1(\by_1) = a_1]} \otimes I_{\reg{Bob}}\nonumber\\
	- &(G_1)^{\bx}_{[g_1(\by_1) = a_1]} \cdot (G_2)^\bx_{[g_2(\by_2) = a_2]} \cdot (G_1)^{\bx}_{[g_1(\by_1) = a_1]} \cdot (G_2)^\bx_{[g_2(\by_2) = a_2]}\otimes I_{\reg{Bob}} \big)\ket{\psi}.\label{eq:what-to-do-with-this}
\end{align}
We can therefore show \Cref{eq:burger-king} is small by
upper-bounding (\Cref{eq:burger-king}$ - $\Cref{eq:what-to-do-with-this}).
There are four terms in this difference;
write $\Delta_i$ for the $i$-th term in \Cref{eq:burger-king} minus the $i$-th term in \Cref{eq:what-to-do-with-this}.
We will bound each $\Delta_i$ one-by-one.

The first term in the difference, $\Delta_1$, is
\begin{multline*}
\E_{\bx, \by_1, \by_2} \sum_{a_1}\sum_{g_2} \bra{\psi}
	(G_2)^{\bx}_{g_2}
	\cdot (G_1)^{\bx}_{[g_1(\by_1) = a_1]} \cdot (G_1)^{\bx}_{[g_1(\by_1) = a_1]} \cdot (G_2)^{\bx}_{g_2} \otimes I_{\reg{Bob}} \ket{\psi}\\
-\E_{\bx, \by_1, \by_2} \sum_{a_1, a_2} \bra{\psi}
	(G_2)^{\bx}_{[g_2(\by_2) = a_2]}
	\cdot (G_1)^{\bx}_{[g_1(\by_1) = a_1]} \cdot (G_1)^{\bx}_{[g_1(\by_1) = a_1]} \cdot (G_2)^{\bx}_{[g_2(\by_2) = a_2]} \otimes I_{\reg{Bob}} \ket{\psi}.
\end{multline*}
The first of these terms is at most~$1$, and so we just have to show that the second term is close to~$1$ as well.
Note that by repeated applications of \Cref{eq:gonna-compare-like-a-pear}, we have that 
\begin{equation*}
(G_2)^{\bx}_{[g_2(\by_2) = a_2]}
	\cdot (G_1)^{\bx}_{[g_1(\by_1) = a_1]} \cdot (G_1)^{\bx}_{[g_1(\by_1) = a_1]} \cdot (G_2)^{\bx}_{[g_2(\by_2) = a_2]} \otimes I_{\reg{Bob}}
	\approx_{\delta} I_{\reg{Alice}} \otimes A^{\bx, \by_1, \by_2}_{a_1, a_2}.
\end{equation*}
But then by \Cref{fact:approx-delta-generalized-game-value}, the expression we want to lower-bound is $O(\delta^{1/2})$-close to
\begin{equation*}
\E_{\bx, \by_1, \by_2} \sum_{a_1, a_2} \bra{\psi} I_{\reg{Alice}} \otimes A^{\bx, \by_1, \by_2}_{a_1, a_2} \ket{\psi},
\end{equation*}
which is exactly~$1$. As a result, $\Delta_1$ is at most $O(\delta^{1/2})$.

The second term in the difference, $\Delta_2$, can be written as
\begin{equation*}
-\E_{\bx, \by_1, \by_2} \sum_{a_1}\sum_{\substack{g_2\neq g_2',\\g_2(\by_2) = g_2'(\by_2)}} \bra{\psi}(
	(G_1)^{\bx}_{[g_1(\by_1) = a_1]} \cdot (G_2)^{\bx}_{g_2} \cdot (G_2)^{\bx}_{g_2'} \cdot (G_1)^{\bx}_{[g_1(\by_1) = a_1]} \otimes I_{\reg{Bob}})\ket{\psi}.
\end{equation*}
This is zero because $G_2$ is a projector.

The third and fourth terms in \Cref{eq:burger-king} are complex conjugates of each other,
as are the third and fourth terms in \Cref{eq:what-to-do-with-this}.
As a result, it suffices to bound the magnitude of $\Delta_4$, and this will serve to bound $\Delta_3$ as well.
We begin by manipulating the fourth term in \Cref{eq:burger-king};
specifically, we will show that it is close to
\begin{equation}\label{eq:cloud}
-\E_{\bx, \by_1, \by_2} \sum_{a_1, g_2} \bra{\psi}
	 (G_1)^{\bx}_{[g_1(\by_1) = a_1]} \cdot (G_2)^{\bx}_{g_2} \cdot (G_1)^{\bx}_{[g_1(\by_1) = a_1]} \otimes (G_2)^{\bx}_{g_2} \ket{\psi}.
\end{equation}
To do so, we take their difference:
\begin{multline*}
\E_{\bx, \by_1, \by_2} \sum_{a_1, g_2} \bra{\psi}
	 ((G_1)^{\bx}_{[g_1(\by_1) = a_1]} \cdot (G_2)^{\bx}_{g_2} \cdot (G_1)^{\bx}_{[g_1(\by_1) = a_1]} \otimes I_{\reg{Bob}})\\
	\cdot (I_{\reg{Alice}} \otimes  (G_2)^{\bx}_{g_2} - (G_2)^{\bx}_{g_2} \otimes I_{\reg{Bob}}) \ket{\psi}.
\end{multline*}
To bound the magnitude, we apply Cauchy-Schwarz:
\begin{multline*}
\E_{\bx, \by_1, \by_2}
	\sqrt{\sum_{a_1, g_1} \Vert  ((G_1)^{\bx}_{[g_1(\by_1) = a_1]} \cdot (G_2)^{\bx}_{g_2} \cdot (G_1)^{\bx}_{[g_1(\by_1) = a_1]} \otimes I_{\reg{Bob}}) \ket{\psi} \Vert^2}\\
	\cdot \sqrt{\sum_{a_1, g_1} \Vert ((G_1)^{\bx}_{[g_1(\by_1) = a_1]} \otimes I_{\reg{Bob}})
	\cdot (I_{\reg{Alice}} \otimes  (G_2)^{\bx}_{g_2} - (G_2)^{\bx}_{g_2} \otimes I_{\reg{Bob}}) \ket{\psi} \Vert^2}.
\end{multline*}
The expression inside the first square root is always at most~$1$.
This allows us to bring the expectation into the second square root by Jensen's inequality.
Because $G_1$ is a POVM, we can bound the resulting expectation by
\begin{equation}\label{eq:majora}
\E_{\bx, \by_1, \by_2} \sum_{g_1} \Vert (I_{\reg{Alice}} \otimes  (G_2)^{\bx}_{g_2} - (G_2)^{\bx}_{g_2} \otimes I_{\reg{Bob}}) \ket{\psi} \Vert^2.
\end{equation}
To bound this, we note that \Cref{eq:gonna-compare-like-a-pear,eq:spend-some-time-reflecting} along with \Cref{fact:triangle-like} imply that
\begin{equation*}
(G_2)^{x}_{[g_2(y_2) = a_2]} \otimes I_{\reg{Bob}}
\consistency_{\delta} I_{\reg{Bob}} \otimes (G_2)^{x}_{[g_2(y_2) = a_2]}.
\end{equation*}
Using the distance properties of $\calG_2$, this implies that
\begin{equation*}
(G_2)^{x}_{g_2} \otimes I_{\reg{Bob}}
\consistency_{\delta + \eps} I_{\reg{Bob}} \otimes (G_2)^{x}_{g_2}.
\end{equation*}
Hence, \Cref{eq:majora} is at most $O((\delta+\eps)^{1/2})$.
A similar argument shows that the fourth term in \Cref{eq:what-to-do-with-this}
is $O(\delta^{1/2})$-close to
\begin{equation}\label{eq:strife}
-\E_{\bx, \by_1, \by_2} \sum_{a_1, a_2} \bra{\psi}
(G_1)^{\bx}_{[g_1(\by_1) = a_1]} \cdot (G_2)^\bx_{[g_2(\by_2) = a_2]} \cdot (G_1)^{\bx}_{[g_1(\by_1) = a_1]} \otimes (G_2)^\bx_{[g_2(\by_2) = a_2]} \ket{\psi}.
\end{equation}

Now, we compute \Cref{eq:cloud} minus \Cref{eq:strife}:
\begin{align*}
&\E_{\bx, \by_1, \by_2} \sum_{a_1}\sum_{\substack{g_2, g_2'\\g_2(\by_2) \neq g_2'(\by_2)}} \bra{\psi}
	 (G_1)^{\bx}_{[g_1(\by_1) = a_1]} \cdot (G_2)^{\bx}_{g_2} \cdot (G_1)^{\bx}_{[g_1(\by_1) = a_1]} \otimes (G_2)^{\bx}_{g_2'} \ket{\psi}\\
=&\E_{\bx, \by_1, \by_2} \sum_{a_1}\sum_{g_2, g_2'} \bra{\psi}
	 (G_1)^{\bx}_{[g_1(\by_1) = a_1]} \cdot (G_2)^{\bx}_{g_2} \cdot (G_1)^{\bx}_{[g_1(\by_1) = a_1]} \otimes (G_2)^{\bx}_{g_2'} \ket{\psi}
	 	\cdot \bone(g_2, g_2', \by_2),
\end{align*}
where $\bone(g_2, g_2', \by_2)$ is the indicator that $g_2 \neq g_2'$ but $g_2(\by_2) = g_2'(\by_2)$.
This is the only part of the expression that depends on $\by_2$, and by our distance assumption it is at most $\eps$ in expectation.
Since the rest of the expression is guaranteed to be positive, we can upper-bound this by
\begin{equation*}
\E_{\bx, \by_1} \sum_{a_1}\sum_{g_2, g_2'} \bra{\psi}
	 (G_1)^{\bx}_{[g_1(\by_1) = a_1]} \cdot (G_2)^{\bx}_{g_2} \cdot (G_1)^{\bx}_{[g_1(\by_1) = a_1]} \otimes (G_2)^{\bx}_{g_2'} \ket{\psi} \cdot \eps.
\end{equation*}
But the remaining part of the expression is at most~$1$, and so in total we can upper-bound it by~$\eps$.
This completes the proof.
\end{proof}

\subsection{Commuting EPR strategies}

In this section, we introduce a class of strategies important for our proof.

\begin{definition}
A strategy $\calS = (\psi, M)$ is called an \emph{EPR strategy} if it satisfies the following two properties.
First, there is an integer~$k$ and powers of two $q_1, \ldots, q_k$ such that
\begin{equation*}
\ket{\psi} = \ket{\mathrm{EPR}_{q_1}} \otimes \ket{\mathrm{EPR}_{q_2}} \otimes \cdots \otimes \ket{\mathrm{EPR}_{q_k}}.
\end{equation*}
Second, for each question~$x$, $M^x$ is a projective measurement. If
for all questions $x$ and answers $a$, $M^x_a$ is a real-valued
matrix, we say that the strategy is \emph{real}

In addition, given a game~$\game$, we say that a real EPR strategy $\calS$ is a \emph{real
  commuting EPR strategy (with respect to~$\game$)}
if for every~$(x_1, x_2)$ in the support of~$\calS$ and every~$a_1, a_2$, $M^{x_1}_{a_1}$ commutes with $M^{x_2}_{a_2}$.
We denote the set of real commuting EPR strategies by $\comstrat{\game}$.
\end{definition}

Real commuting EPR strategies are motivated by the \emph{completeness} cases that arise in this work.
We give a series of transformations which modify games to make them sound against increasingly broader sets of strategies.
Unfortunately, these transformations are not complete for all strategies, in the sense that value-$1$ strategies may be mapped to value-less-than-1 strategies,
but we will be careful to ensure that they \emph{are} complete for all commuting EPR strategies.
For the majority of the paper, the one property of commuting EPR strategies that we will use, not shared by all value-$1$ strategies, is the following.

\begin{fact}\label{fact:heh-heh-heh-gonna-make-anand-prove-this-so-i-can-take-the-day-off}
Let $(\psi, M)$ be a real EPR strategy.
Then $M^x_a \otimes I_{\reg{Bob}} \consistency_0 I_{\reg{Alice}} \otimes M^x_a$ for every distribution on~$x$.
\end{fact}
\begin{proof}
  From the definition of EPR strategies, we know that $\ket{\psi} =
  \ket{\epr_{q_1}} \ot \dots \ot \ket{\epr_{q_k}} \in (\C^{q_1 \cdot
    q_2 \cdot \dots \cdot q_k})^{\ot 2}$. We may choose a basis
  $\{\ket{i} : 1 \leq i \leq q_1 \cdot q_2 \cdot \dots \cdot q_k\}$
  for $\C^{q_1 \cdot \dots \cdot q_k}$, so that
  \[ \ket{\psi} = \sum_{i} \ket{i}_{\reg{Alice}} \ot
    \ket{i}_{\reg{Bob}} . \]
  
Let us denote the components of $M^x_a$ by the notation
$(M^x_a)_{ij}$, so that $M^x_a = \sum_{ij} (M^x_a)_{ij}
\ket{i}\bra{j}$.  Now, for an arbitrary pair $x, a$, we can compute
the post-measurement states from applying $M^x_a$ on Alice's and Bob's systems.
\begin{align*}
  M^x_a \ot I_{\reg{Bob}} \ket{\psi} &= (M^x_a \ot I_{\reg{Bob}}) \sum_{i} \ket{i}_{\reg{Alice}}
                                       \ot \ket{i}_{\reg{Bob}} \\
                                     &= \sum_{ij} (M^x_a)_{ij} \ket{i}
                                       \ot \ket{j} \\
                                     &= \sum_{ij} (M^x_a)_{ji} \ket{i}
                                       \ot \ket{j} \\
                                     &= (I_{\reg{Alice}} \ot M^x_a) \ket{\psi},
\end{align*}
where in going from the second to the third line, we have used the
fact that $M^x_a$ is real and Hermitian, and thus symmetric.
\end{proof}

The following fact is a useful special case.

\begin{fact}\label{fact:the-ol-pauli-swaperoonie}
Let $n > 0$, $q = 2^t$, and $W \in \{X, Z\}$.
Then $\tau^W_u \otimes I \approx_0 I \otimes \tau^W_u$ on the state $\ket{\mathrm{EPR}_q^n}$.
\end{fact}
\begin{proof}
By \Cref{eq:pauli-eigenstates}, we can write 
\begin{equation*}
  \ket{\tau^X_u} = \frac{1}{\sqrt{q}} \sum_{v \in \F_q}
  \omega^{-\tr[uv]} \ket{v},\qquad \ket{\tau^Z_u} = \ket{u}.
\end{equation*}
The second of these self-evidently has real-valued coefficients.
As for the first, $q = 2^t$ implies that $p = 2$.
This means that $\omega = -1$ and $\tr[uv] \in \{0, 1\}$ for all $u, v$.
As a result, it too has real-valued coefficients.
The fact then follows from \Cref{fact:heh-heh-heh-gonna-make-anand-prove-this-so-i-can-take-the-day-off}.
\end{proof}

This property of real commuting strategies is useful for answer
reduction because it allows us to perform oracularization, giving one
prover both questions~$x_1$ and~$x_2$ so that they may simulate the
action of both provers by simultaneously measuring $M^{x_1}$ and $M^{x_2}$.
For more details, see \Cref{part:answer}.

\subsection{Quantum soundness of the classical low-degree test}\label{sec:classical-q-low-deg}

An important tool for quantum protocols
is a version of the Raz-Safra theorem (\Cref{thm:raz-safra})
in which the soundness of the low-degree test is extended to hold even in the case when the provers are allowed to share entanglement.
For the plane-versus-point test, this was first developed by Vidick in~\cite{Vid16}, but for technical reasons he could only show it for the case of three or more quantum provers.
In~\cite{NV18b}, this was improved to hold for the two-prover case, and this is the result we use in this work.
We begin by defining the class of polynomial measurements.

\begin{definition}
Define $\mathrm{PolyMeas}(m, d, q)$ to be the set of POVM measurements whose outcomes correspond to degree-$d$, $\F_q$-valued polynomials.
In other words, $G \in \polymeas{m}{d}{q}$ if $G = \{G_g\}_g$ with outcomes degree-$d$ polynomials $g:\F_q^m \rightarrow \F_q$.
More generally, we let $\simulpolymeas{m}{d}{q}{\ell}$ be the set of measurements $G = \{G_{g_1, \ldots, g_\ell}\}$ outputting~$\ell$ degree-$d$ polynomials $g_i:\F_q^m \rightarrow \F_q$.
\end{definition}

The following theorem establishes the quantum soundness of the classical low-degree test in the $k = 2$ case.

\begin{theorem}[Quantum soundness of the classical low-degree test~{\cite[Theorem 2]{NV18b}}]\label{thm:anand-thomas-classical-low-degree}
There exists a constant $c > 0$ and a function $\delta(\eps) = \mathrm{poly}(\epsilon, dm/q^c)$ such that the following holds.
Suppose Alice and Bob are entangled provers who pass $\game_{\mathrm{Surface}}(m,d,q,2)$ with probability at least $1-\eps$
using the strategy $(\psi, M)$, where~$M$ consists of projective measurements.
Then there exists a POVM measurement $G \in \polymeas{m}{d}{q}$ such that
\begin{equation*}
M^w_b \otimes I_{\reg{Bob}} \consistency_{\delta(\eps)} I_{\reg{Alice}} \otimes G_{[g(w)=b]},
\qquad
G_g \otimes I_{\reg{Bob}} \consistency_{\delta(\eps)} I_{\reg{Alice}} \otimes G_g,
\end{equation*}
where the first is on the uniform distribution over $\F_q^m$.
\end{theorem}

\begin{remark}
The statement of \Cref{thm:anand-thomas-classical-low-degree} is modified from how it appears in \cite[Theorem 2]{NV18b} to better suit our needs.
In this remark, we show how to derive our version from theirs,
which is stated as follows.
\begin{itemize}
\item[$\circ$] There exists a constant $c >0$ and a function $\delta(\eps) = \mathrm{poly}(\eps)$ such that the following holds.
		Suppose $q \geq (dm/\eps)^c$. Then if Alice and Bob pass the surface-versus-point test with probability $1-\eps$, there is a measurement $G \in \polymeas{m}{d}{q}$
		such that
		\begin{equation}\label{eq:what-anand-and-thomas-proved}
		\E_{\bs} \sum_g \sum_{f \neq g|_s} \bra{\psi} M^{\bs}_f \otimes G_g \ket{\psi} \leq \delta(\eps),\qquad
		\sum_g \bra{\psi} G_g \otimes (I - G_g) \ket{\psi} \leq \delta(\eps).
		\end{equation}
\end{itemize}
These are equivalent to the statements
\begin{equation*}
M^s_f \otimes I_{\reg{Bob}} \consistency_{\delta(\eps)} I_{\reg{Alice}} \otimes G_{[g|_s=f]},
\qquad
G_g \otimes I_{\reg{Bob}} \consistency_{\delta(\eps)} I_{\reg{Alice}} \otimes G_g,
\end{equation*}
where the first is on the uniform distribution over surface in $\F_q^m$.
The second of these matches the corresponding statement above.
Next, by \Cref{fact:specialize-the-simeq} we derive
\begin{equation*}
M^s_{[f(w) = b]} \otimes I_{\reg{Bob}} \consistency_{\delta(\eps)} I_{\reg{Alice}} \otimes G_{[g(w)=b]}
\qquad
\text{and}
\qquad
G_{[g(w) = b]} \otimes I_{\reg{Bob}} \consistency_{\delta(\eps)} I_{\reg{Alice}} \otimes G_{[g(w) = b]}.
\end{equation*}
with respect to the distribution $(\bs, \bw)$ from $\game_{\mathrm{Surface}}(m,d,q,2)$.
On top of that, since the strategy passes the test with probability $1-\eps$,
\begin{equation*}
M^w_b \otimes I_{\reg{Bob}} \consistency_{\eps} I_{\reg{Alice}} \otimes M^s_{[f(w) = b]}
\end{equation*}
As a result, if we use \Cref{fact:agreement} to switch these to ``$\approx_\delta$" statements, then
\begin{equation*}
M^w_b \otimes I_{\reg{Bob}}
\approx_{\eps} I_{\reg{Alice}} \otimes M^s_{[f(w) = b]}
\approx_{\delta(\eps)} G_{[g(w)=b]} \otimes  I_{\reg{Bob}} 
\approx_{\delta(\eps)} I_{\reg{Alice}} \otimes G_{[g(w) = b]}.
\end{equation*}
The result now follows from the triangle inequality (\Cref{fact:triangle}) followed by \Cref{fact:almost-agreement} and the fact that~$M$ was assumed to be projective.

Finally, we remove the condition on~$q$ using a trick from~\cite{NV18a}. If $q < (dm/\eps)^c$, then we select $\eps' > \eps$ such that $q = (dm/\eps')^c$.
Alice and Bob also pass the plane-versus-point test with probability $1-\eps'$ because $1-\eps' < 1-\eps$, and so we can apply the theorem with these parameters,
giving a robustness of $\delta(\eps') = \delta(dm/q^{1/c})$.
(In the case when $\eps' > 1$, which is not allowed, this bound trivially still holds because $dm/q^{1/c} > 1$.)
In general, then, we can remove the condition on~$q$ so long as we replace the robustness of $\poly(\eps)$ with $\poly(\eps, dm/q^{1/c})$, which holds in both cases.
\end{remark}

We will use the following proposition about polynomial measurements several times.

\begin{proposition}\label{prop:same-on-point-same-on-subspace}
Let $d > 0$ be an integer.
Consider a distribution $\calD$ on pairs $(\bs, \bu)$, where~$\bs$ is a subspace in $\F_q^m$ and~$\bu$ is a uniformly random point in~$\bs$.
Let $\{M^s_f\}$ be a measurement whose outcomes are degree-$d$ polynomials $f:s \rightarrow \F_q$,
and let $G \in \polymeas{m}{d}{q}$.
Suppose that
\begin{equation*}
M^s_{[f(u) = b]} \otimes I_{\reg{Bob}} \consistency_\delta I_{\reg{Alice}} \otimes G_{[g(u) = b]}
\end{equation*}
with respect to~$\calD$. Then
\begin{equation*}
M^s_{f} \otimes I_{\reg{Bob}} \consistency_{\delta+d/q} I_{\reg{Alice}} \otimes G_{[g|_s = f]}
\end{equation*}
with respect to~$\calD$.
\end{proposition}
\begin{proof}
Suppose the verifier (i)  samples $(\bs, \bu) \sim \calD$,
(ii) gives Alice~$\bs$, who measures with $M^{\bs}$ and returns her outcome~$\boldf : \bs \rightarrow \F_q$,
(iii) receives~$\bg : \F_q^m \rightarrow \F_q$ from Bob, sampled via~$G$,
and (iv) accepts if $\boldf(\bu) = \bg(\bu)$.
Then by assumption, the verifier accepts with probability at least $1-O(\delta)$.

We can use this to bound the probability that $\boldf$ and $\bg$ disagree on the subspace~$\bs$.
By Schwartz-Zippel (\Cref{lem:schwartz-zippel}), conditioned on $\boldf$ and $\bg$ disagreeing,
the probability they disagree on a random point $\bu \sim \bs$  is at least $1-d/q$.
This gives us the inequality $\Pr[\boldf \neq \bg|_{\bs}] \cdot (1-d/q) \leq O(\delta)$.
Now, assume first that $q \geq 2d$. Then this bound implies, via \Cref{fact:agreement}, that
\begin{equation}\label{eq:almost-there-almost-there}
M^s_f \otimes I_{\reg{Bob}} \approx_{\delta + d/q} I_{\reg{Alice}} \otimes G_{[g|_{s} = f]}.
\end{equation}
On the other hand, when $q \geq 2d$, then this bound is also true for trivial reasons.
This is because we can pick $\delta(\cdot)$ such that $\delta(\eps) \geq 1$ in this case.
\end{proof}

\subsection{Quantum soundness of the classical simultaneous low-degree test}\label{sec:q-simultaneous}

We would now like to use \Cref{thm:anand-thomas-classical-low-degree}
to show quantum soundness for the simultaneous classical low-degree test.
This will be done using the same reduction presented in~\Cref{sec:simultaneous-classical}.
The main result is the following.

\begin{theorem}[Quantum soundness of the simultaneous classical
  low-degree test]
  \label{thm:simultaneous-ldt}
There exists a constant $c > 0$ and a function $\delta(\eps) = \mathrm{poly}(\epsilon, d (m + \ell)/q^c)$ such that the following holds.
Suppose Alice and Bob are entangled provers who pass $\game_{\mathrm{Surface}}^\ell(m,d,q,2)$ with probability at least $1-\eps$
using the strategy $(\psi, M)$, where~$M$ consists of projective measurements.
Then there exists a measurement $G \in \simulpolymeas{m}{d}{q}{\ell}$ such that
\begin{equation*}
M^w_{b_1, \ldots, b_\ell} \otimes I_{\reg{Bob}} \consistency_{\delta(\eps)} I_{\reg{Alice}} \otimes G_{[g_1(w), \ldots, g_\ell(w) = b_1, \ldots, b_\ell]},
\qquad
G_{g_1, \ldots, g_\ell} \otimes I_{\reg{Bob}}
\consistency_{\delta(\eps)}
I_{\reg{Alice}} \otimes G_{g_1, \ldots, g_\ell},
\end{equation*}
where the first is on the uniform distribution over $\F_q^m$.
\end{theorem}

\begin{proof}
Suppose Alice and Bob pass $\game_{\mathrm{Surface}}^\ell(m, d, q,2)$ with probability at least $1-\eps$.
We will use them to simulate two provers, ``Combined Alice" and ``Combined Bob", who pass
the single-function low-degree test $\game_{\mathrm{Surface}}(\ell + m, d+1, q,2)$ with probability at least $1-\eps$.
They are specified as follows:
\begin{itemize}
\item[$\circ$] \textbf{Combined Alice:} Given $\bs \subseteq \F_q^{\ell + m}$, draw $\bs' \sim_2 \bs_{\mathrm{proj}}$.  Give it to Alice, who responds with $\boldf_1, \ldots, \boldf_{\ell}: \bs' \rightarrow \F_q$.  Output the function $\mathrm{combine}_{\boldf}|_{\bs}$.
\item[$\circ$] \textbf{Combined Bob:} Given $(\bx, \by) \in \F_q^{\ell + m}$, compute $\by \in \F_q^m$. Give it to Bob, who responds with $\bb_1, \ldots, \bb_\ell \in \F_q$.
			Return $\mathrm{combine}_{\bb}(\bx) \in \F_q$.
\end{itemize}
By \Cref{prop:really-the-same-dist}, $\bs'$ and $\by$ are distributed as the questions in $\game_{\mathrm{Surface}}^\ell(m, d, q,2)$.
Using our assumption on Alice and Bob, this means that $\boldf_1(\by) = \bb_1$, \ldots, $\boldf_\ell(\by) = \bb_\ell$ with probability at least $1-\eps$.
As a result, $(\mathrm{combine}_{\boldf}|_{\bs})(\bx, \by) = \mathrm{combine}_{\bb}(\by)$ with probability at least~$1-\eps$.
By \Cref{prop:sub-subspace}, $\mathrm{combine}_{\boldf}|_{\bs}$ is a degree-$(d+1)$ function on~$\bs$,
and so it is a valid response to subspace queries.
This means Combined Alice and Bob pass~$\game_{\mathrm{Surface}}(\ell + m, d+1, q, 2)$ with probability at least~$1-\eps$.

Thus, we can apply \Cref{thm:anand-thomas-classical-low-degree}.
It gives a measurement $G \in \polymeas{\ell+m}{d+1}{q}$ such that
\begin{equation}\label{eq:g-meas-equals-combine-meas}
M^y_{[\mathrm{combine}_b(x)=\nu]} \otimes I_{\reg{Bob}}
\consistency_{\delta(\eps)} I_{\reg{Alice}} \otimes G_{[g(x, y) = \nu]},
\qquad
G_g \otimes I_{\reg{Bob}} \consistency_{\delta(\eps)} I_{\reg{Alice}} \otimes G_g,
\end{equation}
where $\delta(\eps) = \poly(\eps, (d+1)(\ell+m)/q^c)$.
This means that if we give Alice~$\by$ and she returns~$\bb$, and Bob simply returns~$\bg$,
then $\mathrm{combine}_{\bb}(\bx) = g(\bx, \by)$ with probability at least~$1-\delta(\eps)$.

We would like to show that~$\bg$ is exactly linear in~$x$ with high probability, over the randomness in the measurement~$G$.
Let us condition on a $\bg$ which is not exactly linear.
By \Cref{prop:exactly-linear-prop}, the probability that $\bg|_{\by}$ is not exactly linear is at least $1-(d+1)/q$.
On the other hand, because $\mathrm{combine}_{\bb}(\bx)$ is always exactly linear by construction,
the probability that $\bg|_{\by}(\bx) = \mathrm{combine}_{\bb}(\bx)$ is at most $(d+1)/q$ by Schwartz-Zippel (\Cref{lem:schwartz-zippel}).
As a result, the probability that $\bg(\bx, \by) = \mathrm{combine}_{\bb}(\bx)$ is at most $(d+1)/q + (d+1)/q$.
Thus, if we write $\mu_{\mathrm{linear}}$ for the probability that $\bg$ is exactly linear, we have equality at most $\mu_{\mathrm{linear}} + 2(d+1)/q$ fraction of the time.
Rearranging, $\bg$ is exactly linear with probability
\begin{equation*}
\mu_{\mathrm{linear}} \geq 1 - 2\delta(\eps) - 2(d+1)/q.
\end{equation*}
Define a new measurement $\{H_{g_1, \ldots, g_\ell}\} \in \simulpolymeas{m}{d}{q}{\ell}$ operationally as follows: first, measure~$G$ and receive~$\bg$.
If it is exactly linear, it can be written as $\sum_i x_i \cdot \bg_i(y)$, and so output $\bg_1, \ldots, \bg_\ell$.
If~$\bg$ is \emph{not} exactly linear, output any arbitrary degree-$d$ polynomials instead.
When~$\bg$ is exactly linear, we have $\mathrm{combine}_{\bg_1, \ldots, \bg_\ell}(x, y) = \bg(x, y)$.
Since this happens with probability at least $1-\delta(\eps)$,
we can replace~$G$ with~$H$ in \Cref{eq:g-meas-equals-combine-meas}, yielding
\begin{equation}\label{eq:no-name-is-a-bad-name}
M^y_{[\mathrm{combine}_b(x)=\nu]} \otimes I_{\reg{Bob}}
\consistency_{\delta(\eps)} I_{\reg{Alice}} \otimes H_{[\mathrm{combine}_g(x, y) = \nu]},
\end{equation}
On the other hand, because~$H$ is just~$G$ with data processing applied to its output,
we can apply \Cref{fact:specialize-the-simeq} to \Cref{eq:g-meas-equals-combine-meas}. This produces the equation
\begin{equation*}
H_{g_1, \ldots, g_\ell} \otimes I_{\reg{Bob}} \consistency_{\delta(\eps)} I_{\reg{Alice}} \otimes H_{g_1, \ldots, g_\ell}.
\end{equation*}

Consider $\bg_1, \ldots, \bg_\ell$ drawn by Bob using~$H$.
For any fixed~$b$ and~$y$, if it is not the case that $g_1(y) = b_1$, \ldots, $g_\ell(y) = b_\ell$,
then the probability that $\mathrm{combine}_{\bg}(\bx, y) = \mathrm{combine}_{b}(\bx)$
over a random~$\bx$ is at most $1/q$ by Schwartz-Zippel, since both are exactly linear functions.
Thus, if Alice draws~$\bb_1, \ldots, \bb_\ell$ given~$\by$,
and we write~$\eta$ for the probability that $\bg_1(\by) = \bb_1$, \ldots, $\bg_\ell(\by) = \bb_\ell$,
then the probability that $\mathrm{combine}_{\bg}(\bx, \by) = \mathrm{combine}_{\bb}(\bx)$ is at most
$\eta +  (1-\eta) \cdot 1/q$.
Combined with \Cref{eq:no-name-is-a-bad-name}, this implies that
\begin{equation*}
\Pr[\bg_1(\by) = \bb_1, \ldots, \bg_\ell(\by) = \bb_\ell] \geq 1- \delta(\eps) - 1/q.
\end{equation*}
Or, equivalently,
\begin{equation*}
M^y_{b_1, \ldots, b_\ell} \otimes I_{\reg{Bob}}
\consistency_{\delta(\eps)} I_{\reg{Alice}} \otimes G_{[g_1(y), \ldots, g_\ell(y) = b_1, \ldots, b_\ell]}.\qedhere
\end{equation*}
\end{proof}

\subsection{Self-testing}

The games presented in \Cref{sec:classical-q-low-deg,sec:q-simultaneous} might be referred to as ``measurement testers":
if a strategy passes them with high probability, then we can extract some property on its measurements.
In this section, we will introduce a significantly stronger notion of testing called \emph{self-testing}.
A self-tester is a game in which if a prover passes with high probability,
then not only do we do exactly which measurements the prover must be performing,
we also know which exactly state it must be performing them on (up to local isometry).
(We note that some works use ``self-testing" to refer both to ``measurement testing" and what we refer to as ``self-testing"~\cite{NV18a}.
In this work, we will reserve the term exclusively for the latter.) We begin with a definition.

\begin{definition}
We say that $\mathcal{S} = (\psi, M)$ is a \emph{partial strategy} for $\game$
if $M$ contains the POVM $M^x$ for only a subset of the questions in~$\game$.
We call this set of questions $\mathcal{S}$'s \emph{question set}.
A strategy $\mathcal{S}' = (\psi, M')$ \emph{extends} $\mathcal{S}$ if $(M')^x = M^x$ for every~$x$ in $\mathcal{S}$'s question set.
\end{definition}

Next, we define self-testing.

\begin{definition}[Self-testing]
Let $\mathcal{S} = (\psi, G)$ be a partial strategy and $\mathcal{D}$ be a distribution over its question set.
A game $\game$ is a \emph{self-test for $\mathcal{S}$ over $\mathcal{D}$ with robustness $\delta(\epsilon)$}
if it satisfies the following two conditions.
\begin{itemize}
\itemsep -.5pt
\item[$\circ$] \textbf{Completeness:} There exists a (full) strategy $\mathcal{S}_{\mathrm{full}}$ consistent with $\mathcal{S}$ which passes~$\game$ with probability~$1$.
\item[$\circ$] \textbf{Soundness:} Let $\overline{\mathcal{S}} = (\overline{\psi}, \overline{M})$ be a strategy which passes~$\game$ with probability $1-\eps$.
	Then there exists a local isometry $\phi = \phi_{\mathrm{local}} \otimes \phi_{\mathrm{local}}$ and a state $\ket{\mathrm{aux}}$ such that
	\begin{equation*}
		\Vert \phi \ket{\overline{\psi}} - \ket{\psi}\ket{\mathrm{aux}} \Vert^2 \leq \delta(\epsilon).
	\end{equation*}
	\ignore{and
	\begin{equation*}
		\E_{(\bx, \bx') \sim\game} \sum_{a, a'} \Vert \phi (M_a^{\bx} \otimes M_{a'}^{\bx'} \ket{\psi'})
				- (M^x_a \otimes M^x_{a'} \ket{\psi}) \ket{\mathrm{aux}} \Vert^2 \leq \delta(\epsilon).
	\end{equation*}}
	Furthermore, if we define the new matrices $M_a^x := \phi_{\mathrm{local}} \cdot \overline{M}_a^x \cdot (\phi_{\mathrm{local}})^\dagger$, then
	\begin{equation}\label{eq:cite-once-then-never-again}
	M_a^x \otimes I_{\reg{Bob}} \approx_{\delta(\epsilon)} (G_a^x \otimes I_{\mathrm{aux}}) \otimes I_{\reg{Bob}},
	\end{equation}
	on states $\ket{\psi}\ket{\mathrm{aux}}$ and $\ket{\psi'}$ and distribution $\bx \sim \calD$.
\end{itemize}
\end{definition}

We note that this definition of self-testing differs in several key places from the one given in \cite[Definition~$2.5$]{NV18a}.
We will explain these differences in more detail when we cite the quantum low-degree test in \Cref{sec:self-test-pauli}.

%%% Local Variables:
%%% mode: latex
%%% TeX-master: "main"
%%% End:
%%%---------- close: prelims-quantum

\part{Implementing the registers}

\label{part:stack}

%%%%%%%%%%% jump to register-overview
%%%---------- open: register-overview
%!TEX root = main.tex

